\documentclass[11pt]{article}
\usepackage[margin=1in]{geometry}
\usepackage{amsmath,amssymb,amsthm,booktabs,hyperref}
\hypersetup{hypertexnames=false,colorlinks=true,urlcolor=blue,citecolor=blue,linkcolor=blue}
\setlength{\emergencystretch}{1em}
\newtheorem{theorem}{Theorem}
\newtheorem{lemma}[theorem]{Lemma}
\newtheorem{proposition}[theorem]{Proposition}
\newtheorem{remark}[theorem]{Remark}
\newcommand{\PP}{\mathcal P_N}
\newcommand{\li}{\operatorname{li}}
\newcommand{\om}{\omega}
\newcommand{\Th}{U_NN/L^2}
\title{Exact and finite presieving in explicit Chen bounds:\\a GRH threshold and quantitative distribution interfaces}
\author{Research manuscript}
\date{4 October 2026}
\begin{document}
\maketitle
\begin{abstract}
We study the gap between explicit Chen bounds and the verified binary
Goldbach range. We combine exact weights for the smallest primes,
finite Rosser weights for medium primes, and Rosser weights for larger
primes; truncated Brun weights provide an alternative construction.
Finite first-failure masses use exponents selected per degree and
nonnegative tilted squares on a band supplied by the preceding prefix.
Positive prefix blocks retain preceding checks. A six-state count
bounds the union of accepted supports; an integer roughness bound
and a finite analytic interval reduce the GRH threshold.
The actual Rosser support is also smaller than its defining level.
A product condition uniform at every
upper endpoint is used throughout; the problematic fixed-endpoint-only
application in earlier candidates is not needed. The composite weights
retain coefficients of absolute value at most one and have support
below $Q_*H_0D$, where $H_0$ bounds the medium-prime support and
$Q_*$ contains only the exactly presieved primes,
avoiding a full primorial cost through $T$. Using published
explicit sieve constants and the GRH AP and bilinear estimates of
Bordignon--Starichkova, we obtain $N\ge\exp(24274)$ with normalized
surplus exceeding $10^{-5}$, and the stronger surplus $0.0181$
for $N\ge\exp(31000)$. Quantitative signed distribution inputs at levels
$3/4$, $9/10$, $19/20$, and $99/100$ give parameterized thresholds
without RH or GRH, retaining unspecified constants and starting points.
A quantified distribution input starting at $4\cdot10^{18}$ is
withdrawn after an integrality obstruction. This unreviewed manuscript
states its numerical certificates and analytic dependencies separately.
\end{abstract}

\section{Statements, normalization, and comparison}
Put $L=\log N$, and let $\Omega$ count prime factors with multiplicity.
Write $\pi_2(N)$ for the number of primes $p<N$ for which
$N-p\ge2$ and $\Omega(N-p)\le2$. Define
\[
C_\infty=\prod_{p>2}\left(1-\frac1{(p-1)^2}\right),\quad
C_N=C_\infty\prod_{\substack{p>2\\p\mid N}}\frac{p-1}{p-2},\quad
U_N=2e^{-\gamma}C_N.
\]
Since the product over \emph{all} integers $n\ge3$ of
$1-(n-1)^{-2}$ equals $1/2$, $C_N>1/2$ and $U_N>e^{-\gamma}$.
The finite certificate also gives $C_\infty>0.66$, hence the stronger
uniform lower bound $U_N>U_*:=1.32e^{-\gamma}$ used below.
This normalization follows the actual sieve product; it is not the
$2e^{\gamma}C_N$ printed in the accessible GRH preprint.

\begin{theorem}[GRH threshold]\label{thm:grh}
Assume GRH for Dirichlet $L$-functions, including RH for $\zeta$.
Every even integer $N\ge e^{24274}$ satisfies
\[
\pi_2(N)>10^{-5}\,\frac{U_NN}{\log^2N}>0.
\]
For $N\ge e^{31000}$ the coefficient $10^{-5}$ can be replaced by $0.0181$.
\end{theorem}
The published GRH result \cite{BSpub} has threshold $e^{e^{14}}$;
the older preprint \cite{BS} has $e^{e^{15.85}}$.
Here $\log(24274)=10.097161097\ldots$. The new threshold has about
$10543$ decimal digits, compared with about $522284$ for the published
threshold. Theorem \ref{thm:grh} is a result of this unreviewed manuscript,
not a published record. The stronger earlier candidates $e^{8500}$ and
$e^{15000}$ remain unaudited at their original sieve interface and are
not asserted by this theorem.
We do not identify the internal unconditional manuscript at
$e^{e^{24.9}}$ with an independently audited published result.

Let $B=4\cdot10^{18}$. The following \emph{rejected candidate}, which
was previously proposed as an additional hypothesis, is retained for audit:
\begin{equation}\label{eq:QEH}
\mathcal E_{9/10}(x):=
\sum_{d\le x^{9/10}}\mu^2(d)
\max_{(a,d)=1}\left|\pi(x;d,a)-\frac{\pi(x)}{\varphi(d)}\right|
\le\frac{x}{10\log^2x}\qquad(x\ge B).
\tag{QEH$_B$}
\end{equation}
\begin{remark}[Withdrawn full-range claim]\label{thm:bridge}
The former full-range claim based on \eqref{eq:QEH} is withdrawn.
The hypothesis itself is false at $B$, as proved below. Positive
sieve-budget arithmetic under an inconsistent premise is not a
meaningful conditional bridge.
\end{remark}
Ordinary Elliott--Halberstam is a standard asymptotic conjecture,
with constants depending on the distribution level and logarithmic
power. It does not provide this particular finite starting point
or constant. Conversely, a single-level, single-logarithmic-power
finite condition would not imply all of ordinary EH. They must not
be described as two logically ordered versions of the same conjecture.
The original published Chen theorem assumes GRH alone, and the GRH
theorem above keeps that same conjectural assumption. RH for the zeta
function is included in GRH here. EH and GEH are separate recognized
distribution conjectures: for example, the GPY, Maynard, and Polymath
conditional prime-gap results in \cite{Polymath} explicitly invoke
them. GRH is not known to imply either conjecture. A hypothesis with
prescribed numerical constants and a prescribed finite starting point
must be identified separately; the standard asymptotic notation does
not supply those data.
Thus a distribution estimate with unspecified constants can be a
consequence of standard EH, while the same estimate with prescribed
constants and a prescribed starting point is additional quantitative
input. The four distribution routes below use EH-style input in place
of GRH, rather than requiring GRH together with EH.

\begin{proposition}[Integrality obstruction]\label{prop:integrality}
For $2\le H<x$, put $P=\pi(x)$. Then
\[
\sum_{d\le H}\mu^2(d)\max_{(a,d)=1}
\left|\pi(x;d,a)-\frac{P}{\varphi(d)}\right|
\ge\frac{\pi(H)}2-\sqrt x.
\]
Consequently \eqref{eq:QEH} fails at $x=B$ unconditionally.
\end{proposition}
\begin{proof}
For $x<3$ the lower bound is negative and the assertion is immediate.
We may therefore assume $x\ge3$, so $P-1\ge1$.
Restrict the sum to prime moduli $q\le H$, and put $m=q-1$.
The integer counts over reduced classes sum to $P-1$, because
the prime $q$ itself is omitted. Write $P=km+r$, $0\le r<m$.
If $r=0$, some class has count at most $k-1$, so its error is at
least one. If $2\le r<m$, there are classes with count at most $k$
and at least $k+1$, so the maximum error is at least
$\max(r/m,1-r/m)\ge1/2$. The remaining case $r=1$ requires
$q-1\mid P-1$. There are at most $\tau(P-1)\le2\sqrt{P-1}<2\sqrt x$
such exceptional prime moduli. This proves the lower bound.
At $B$, take $H=B^{9/10}$ and $J=\lfloor H\rfloor$.
By \cite[Lemma 3.1]{JS}, $\pi(J)>J/\log J$ for $J\ge71$.
The Arb certificate gives
\[
\frac{J}{2\log J}-\sqrt B>7.15819\cdot10^{14},\qquad
\frac{B}{10\log^2B}<2.18026\cdot10^{14}.
\]
The lower bound already exceeds the proposed upper bound by a
factor greater than $3.28$. The earlier candidate with level $3/4$
and budget $B/\log^4B$ also fails this check, by a factor greater than
$1.169$. Both are rejected, independently of RH or GRH.
\end{proof}

\begin{proposition}[Distribution constants and onsets]\label{prop:tradeoff}
Suppose $0<\theta<1$ and
$\mathcal E_\theta(x)\le Cx/\log^4x$ for every $x\ge X$,
where $C>0$, $X>1$, and the complete square-free modulus sum is
defined in \eqref{eq:EHinput}.
For $y\ge\log X$ with $e^{\theta y}\ge72$, necessarily
\begin{equation}\label{eq:integralitytradeoff}
\begin{aligned}
C&>\mathfrak I_\theta(y),\\
\mathfrak I_\theta(y)
&=\frac{y^3}{2\theta}
 \bigl(e^{-(1-\theta)y}-e^{-y}\bigr)-y^4e^{-y/2}.
\end{aligned}
\end{equation}
Let
$R=\max\{\log X,A,\sqrt{bC}\}$, where $A>0$ and $b>0$.
If $\mathfrak I_\theta(M)>M^2/b$ and $e^{\theta M}\ge72$,
then $R>M$. The following are strict, certified necessary bounds,
with $B=4\cdot10^{18}$:
\[
\begin{array}{c|r|r|r|r}
\theta & C\text{ if }X\le B & \log X\text{ if }C=1 & b & R\\ \hline
3/4    & >1.1697    & >43   &20& >22\\
9/10   & >602.35    & >143  &10& >57\\
19/20  & >4858.17   & >336  &10& >130\\
99/100 & >25860.94  & >2246 &10& >834
\end{array}
\]
These are necessary constraints, not sufficient choices of $C,X$.
The witnesses are not asserted to give optimal lower limits.
\end{proposition}
\begin{proof}
Put $x=e^y$, $H=x^\theta$, and $J=\lfloor H\rfloor\ge71$.
Proposition \ref{prop:integrality} and the prime-count bound in
\cite[Lemma 3.1]{JS} give
\[
\mathcal E_\theta(x)>
\frac{J}{2\log J}-\sqrt x>
\frac{H-1}{2\log H}-\sqrt x.
\]
Multiplying by $y^4/e^y$ gives \eqref{eq:integralitytradeoff}.
For the endpoint constraint, suppose $R\le M$. Then $e^M\ge X$,
so $C>\mathfrak I_\theta(M)>M^2/b\ge R^2/b\ge C$, a contradiction.
The first numerical column evaluates \eqref{eq:integralitytradeoff}
at $y=\log B$. The second uses the witnesses $y=43,143,336,2246$,
where $\mathfrak I_\theta(y)>1$. The last uses
$M=22,57,130,834$, respectively. The 192-bit Arb certificate checks
each strict inequality, including its domain condition, and compares
the closed formula with the separately evaluated normalized
prime-count lower bound.
\end{proof}
In particular, the level $99/100$ thresholds in Theorem
\ref{prop:EH} are always greater than $e^{834}$ under their present
budget, even though their displayed analytic starts are $141$ and $142$.
If $C=1$, the same distribution input instead requires $X>e^{2246}$.
The other fixed analytic starts still dominate the corresponding
last-column bounds. A higher distribution level alone does not
justify a lower effective threshold.

\section{Product estimates and an unused fixed-endpoint candidate}
For $t>2$, use strict upper endpoints and put
\[
D(t)=\log t\prod_{2<p<t}\left(1-\frac1{p-1}\right),\qquad
D_\infty=2e^{-\gamma}C_\infty.
\]
The identity
\[
\left(1-\frac1{p-1}\right)
=\left(1-\frac1p\right)\left(1-\frac1{(p-1)^2}\right)
\]
explains the normalization and avoids bounding two cancelling terms
independently. We use the following two primary inputs:
\begin{align}
\sum_{p<t}\frac1p&\ge\log\log t+M-\frac{2.964\cdot10^{-6}}{\log t}
&& (t\ge2),\label{eq:recip}\\
\left|e^\gamma\log t\prod_{p<t}(1-1/p)-1\right|
&\le\delta(t):=\frac{3\log t}{8\pi\sqrt t}
&& (t\ge23.8,\ \mathrm{RH}).\label{eq:mertens}
\end{align}
These are respectively \cite[Lemma 3.2]{JS} and
\cite[Corollary 1.3]{LN}. The strict endpoint in the second follows
by a one-sided limit of the published non-strict version.

Let $T=10^6$ and set
\[
r_T=\frac{2.964\cdot10^{-6}}{\log T}+\frac1{T-1}+\frac{0.54}{T}
<1.754543\cdot10^{-6}.
\]
\begin{lemma}\label{lem:product}
The certificate described below proves $D(u)<D_\infty$ for
$11<u\le T$. For $u>T$, $D(u)/D_\infty<e^{r_T}$.
Consequently, under RH, for every $z>T$ and $1<u<z$,
\begin{equation}\label{eq:product}
\prod_{\substack{u\le p<z\\p>11}}
\left(1-\frac1{p-1}\right)^{-1}
<\frac{e^{r_T}}{1-\delta(z)}\frac{\log z}{\log u}.
\end{equation}
In particular the factor preceding $\log z/\log u$ is less than
$1+2\cdot10^{-6}$ for $z\ge e^{40}$, and less than $1+1/700$
for $z\ge B^{1/3}$.
\end{lemma}
\begin{proof}
Between consecutive primes $D(u)$ is increasing, and its downward
jumps occur just after a prime. It is enough to check $D(p)$ for
prime $11<p\le T$ and the endpoint $T$. Let
$C_T=\prod_{2<p\le T}(1-(p-1)^{-2})$. Since
\[
\log\frac{C_T}{C_\infty}
\le\sum_{p>T}\frac1{p(p-2)}<\frac1{T-1},
\]
we may compare every peak with the rigorous lower bound
$2e^{-\gamma}C_Te^{-1/(T-1)}$.
The script \path{conditional_chen_product_certificate.py}
uses a complete Eratosthenes sieve and 192-bit Arb arithmetic.
It checks 78493 peaks; the largest occurs at $p=463181$. Its ball
lies below $0.741284513$, whereas the lower bound for
$D_\infty$ exceeds $0.741307533$. These are strict ball comparisons,
not rounded floating-point decisions.

For $p>T$, put
$r(p)=-\log(1-1/(p-1))-1/p$. The factorization above gives
\[
0<r(p)\le\frac1{p(p-2)}+\frac{0.54}{p^2}.
\]
Writing $E(u)=\sum_{p<u}1/p-\log\log u-M$ gives the exact identity
\[
\log\frac{D(u)}{D_\infty}=-E(u)+\sum_{p\ge u}r(p).
\]
Use \eqref{eq:recip} and bound the tails by sums over all integers
larger than $T$. This proves the claimed $r_T$ bound.
Also \eqref{eq:mertens} and $C(z)\ge C_\infty$ give
$D(z)/D_\infty\ge1-\delta(z)$. For $u>11$, the product in
\eqref{eq:product} equals $D(u)\log z/(D(z)\log u)$.
For $u\le11$, take a limit $u'\downarrow11$ from above; no factor
changes and $\log u'\ge\log u$. Finally $\delta(z)$ decreases
for $z>e^2$; the two endpoint comparisons are certified in the scripts.
\end{proof}

Deleting primes dividing $N$, or reducing any local density below
$1/(p-1)$, only decreases the left side. We take $\PP$ to be the
odd primes not dividing $N$ and presieve its primes at most 11.
The actual $Q_N$ divides $1155$, hence $\log Q_N<7.1$.
We occasionally use the weaker bound $\log Q_N<8$.

\begin{remark}[Endpoint quantifiers]\label{rem:endpoints}
Lemma \ref{lem:product} is a condition at the specified upper endpoint
$z$, for every smaller $u$. It is not a uniform condition for all
intermediate upper endpoints: at $u=17,v=19$ the normalized product
equals
$\frac{16}{15}\frac{\log17}{\log19}>1.0263$.
The wording of \cite[Theorem 6]{BJS}, also reproduced in
\cite[Lemma 2.2]{DJ}, is a fixed-upper-endpoint condition. However,
the proof of BJS Lemma 11 applies an induction at lower upper endpoints,
through the recursion $T_n(D,z)=\sum_{p<z}g(p)T_{n-1}(D/p,p)$
in its relevant ranges. The product hypothesis at such an endpoint $p$
does not follow simply by dividing two inequalities at $z$.
We have not verified an alternative argument that closes this gap for
the present product. This is a pending audit of \emph{our application},
not a claim that the published theorem is false.
Neither a positive budget nor a finer parameter scan settles this
interface. The earlier candidates using only this lemma remain pending.
Theorem \ref{thm:grh} instead uses the uniform product estimate and
finite presieving proved next, so its induction has the required
premise at every intermediate endpoint.
\end{remark}

\section{Uniform products and composite sieve weights}
For an integer $T\ge10^4$, put
\[
b_T=\max\left\{\frac2{\sqrt T\log T},
\frac{6.836\cdot10^{-6}}{\log(10^{12})}\right\},\quad
K_T=\exp\left(b_T+\frac{2.964\cdot10^{-6}}{\log T}
+\frac1{T-1}+\frac{0.54}{T}\right).
\]
\begin{lemma}[Uniform product condition]\label{lem:uniform}
For every $1<u<v$, unconditionally,
\[
\prod_{\substack{u\le p<v\\p>T}}
\left(1-\frac1{p-1}\right)^{-1}
\le K_T\frac{\log v}{\log u}.
\]
Deleting primes or decreasing local densities preserves this bound.
\end{lemma}
\begin{proof}
BJS Lemma 16 gives $E(v)\le2/(\sqrt v\log v)$ for $v\le10^{12}$;
JS Lemma 3.2 gives $E(v)\le6.836\cdot10^{-6}/\log v$ for
$v>10^{12}$. Strict endpoints follow by limits, so $E(v)\le b_T$
for $v>T$. For $T<u<v$, use
\[
\log\frac{D(u)}{D(v)}
=E(v)-E(u)+\sum_{p\ge u}r(p)-\sum_{p\ge v}r(p)
\le b_T+\frac{2.964\cdot10^{-6}}{\log T}
+\frac1{T-1}+\frac{0.54}{T}.
\]
This proves the bound. For $u\le T<v$, take $u'\downarrow T$
from above and use $\log u'\ge\log u$. For $v\le T$ the product
is empty. The argument includes all intermediate upper endpoints.
\end{proof}

\begin{lemma}[Finite ordered products with an analytic tail]\label{lem:ordered}
Let $10^4\le T<S\le10^{12}$ be even integers, put $d(t)=\log D(t)$,
and let $\lambda_Q$ be the explicit lower bound
\[
\lambda_Q=\log2-\gamma+
\sum_{2<p\le Q}\log(1-(p-1)^{-2})+\log(1-1/Q)
\le\log D_\infty\qquad(Q\ge3).
\]
Suppose a finite certificate supplies $A_T,M_T$ such that
\[
A_T\ge\sup_{T\le u<v\le S}(d(u)-d(v)),\qquad
M_T\ge\sup_{T\le u\le S}(d(u)-\lambda_Q).
\]
Writing $b_S,K_S$ for the original bounds above, set
\[
B_T=\max\{A_T,M_T+b_S,\log K_S\}.
\]
Then Lemma \ref{lem:uniform} holds with $K_T$ replaced by $e^{B_T}$.
It remains unconditional and uniform at every intermediate endpoint.
\end{lemma}
\begin{proof}
The prime-only tail of $C_\infty$ is at least the product over all
integers $n>Q$:
\[
\prod_{n>Q}(1-(n-1)^{-2})=(Q-1)/Q.
\]
This proves the asserted lower bound $\lambda_Q$.
For $T<u<v\le S$, use the bound $A_T$. For $T<u\le S<v$,
the identity for $D(v)/D_\infty$ gives
$d(v)-\log D_\infty=-E(v)+\sum_{p\ge v}r(p)\ge-b_S$,
so $d(u)-d(v)\le M_T+b_S$.
For $S\le u<v$, apply the original uniform bound with cutoff $S$;
$S$ is composite, so it introduces no missing boundary prime.
The cases $u\le T$ and $v\le T$ follow as in Lemma \ref{lem:uniform}.
Finally, $D(u)/D(v)$ is the product in that lemma multiplied by
$\log u/\log v$.

The function $d(t)$ increases strictly between primes and has a
negative jump at each prime. Thus its ordered drawdown on $[T,S]$
is bounded by checking every left limit as a potential earlier
maximum and every right limit as a potential later minimum, together
with the initial endpoint and the final endpoint.
For adjacent finite intervals with peak, trough and drawdown
enclosures $(P_1,Q_1,A_1)$ and $(P_2,Q_2,A_2)$, the combined
drawdown is bounded by
$\max\{A_1,A_2,P_1-Q_2\}$.
This gives a finite, complete certificate for the continuum of ordered
pairs, rather than a sample of individual endpoints.
\end{proof}
In the applications below take $S=10^8$, $Q=10^7$, and use
$\widehat K_T=e^{B_T}$ from Lemma \ref{lem:ordered}.
The complete integer sieve generates all $5761454$ odd primes through
$S$ and checks $5760226$ prime jumps above the smallest cutoff.
Strict Arb comparisons certify the following rational upper bounds:
\[
\begin{array}{r|r|c}
T&10^{12}\log\widehat K_T&\widehat K_T-1\ \text{(upper bound)}\\\hline
10000&987086120&0.000987573450\\
22000&582114805&0.000582284267\\
70000&291898888&0.000291941495\\
5000000&24260241&0.000024260536
\end{array}
\]
Here the middle column is a certified upper bound, not an asserted
exact maximum. Finite and infinite ranges are joined by the three
cases in the proof; no conjectural prime-product estimate is used.

Let $g(p)=1/(p-1)$ on the odd primes not dividing $N$. Choose
$2\le w_0<T<z$ and split $P_N(z)$ into the exact part $P_E$ with
$p\le w_0$, the medium-prime part $P_0$ with $w_0<p\le T$,
and the large-prime part $P_1$ with $T<p<z$. Put
\[
\begin{aligned}
Q_0&=P_E,& Q_*&=\prod_{2<p\le w_0}p,\\
V_*^{(0)}&=\prod_{2<p\le w_0}\left(1-\frac1{p-1}\right),&
W_*^{(0)}&=\sum_{2<p\le w_0}\frac1{p-1}.
\end{aligned}
\]
Then $Q_0\mid Q_*$. Exact inclusion--exclusion on $Q_0$ has
divisor sum $1_{(n,Q_0)=1}$ and support at most $Q_*$.
Fix an even integer $k\ge2$.
Define Brun coefficients $b^+(d)=\mu(d)$ for $d\mid P_0$,
$\om(d)\le k$, and zero otherwise; define $b^-(d)$ with cutoff
$k-1$. The binomial identity
$\sum_{j=0}^k(-1)^j\binom rj=\binom{r-1}k$ for $r\ge1$
proves the pointwise upper bound, with the analogous negative identity
for $k-1$ proving the lower bound. Put
\[
V_0=\prod_{p\mid P_0}(1-g(p)),\quad
B^\pm=\sum_{d\mid P_0}b^\pm(d)g(d),\quad
\Delta=B^+-B^-.
\]
Then $B^-\le V_0\le B^+$ and
\[
0\le\Delta=\sum_{\substack{d\mid P_0\\\om(d)=k}}g(d)
\le W_{w_0,T}^k/k!,\qquad B^+\le(1+\eta)V_0,
\quad\Delta\le\eta V_0,
\]
provided
\begin{equation}\label{eq:Bruneta}
\begin{aligned}
W_{w_0,T}&=\log\log T+0.262-\frac12+c_0+b_T-W_*^{(0)},\\
c_0&=\sum_{2<p\le1000}\frac1{p(p-1)}+\frac1{1000},\\
\eta&\ge\frac{V_*^{(0)}\log T\,e^{b_T}}{U_*}
\frac{W_{w_0,T}^k}{k!}.
\end{aligned}
\end{equation}
The inequalities $B^-\le V_0\le B^+$ follow by averaging the
pointwise Brun inequalities over independent prime-divisibility
indicators with probabilities $g(p)$. Their difference is the
nonnegative elementary symmetric sum of degree $k$ displayed above.
Indeed $M<0.262$, and the integer telescoping tail bounds the
correction from $1/p$ to $1/(p-1)$ by $c_0$. The exact identity
for $D(T+)$ gives the complete odd-prime product through $T$
at least $U_*e^{-b_T}/\log T$. Dividing this product by
$V_*^{(0)}$ gives a lower bound for the complete medium-prime product,
hence
\[
V_0\ge\frac{U_*e^{-b_T}}{V_*^{(0)}\log T}.
\]
The medium density sum is at most $W_{w_0,T}$. These two bounds
remain valid when primes dividing $N$ are omitted, including omitted
primes in the exact part: the bounds use the complete medium-prime
product and sum, not $V(P_E)$ as a substitute for $V_*^{(0)}$.
The bound $M<0.262$ is independently certified by the partial sum
$\gamma+\sum_{p\le1000}(\log(1-1/p)+1/p)$, whose omitted terms
are negative. The bound $C_\infty>0.66$ is certified by the finite
product through $10^4$, times $e^{-1/9999}$.

There is a second medium-prime construction that uses the divisor
product rather than its number of prime factors to restrict support.
It needs no dimension-one product premise on the medium primes.
\begin{lemma}[Finite Rosser presieving]\label{lem:rankin}
Let $D_0\ge T^2$ and $\sigma>0$. On $P_0$, take $b^\pm(d)=\mu(d)$
if, for $d=p_1\cdots p_n$ with $p_1>\cdots>p_n$, every odd
prefix $m\le n$ (upper sign), or every even prefix (lower sign),
satisfies $p_1\cdots p_mp_m^2<D_0$; otherwise take zero.
These weights have support $d<D_0$, coefficients in $\{\mu(d),0\}$,
and pointwise divisor sums bracketing $1_{(n,P_0)=1}$.
With $B^\pm,V_0,\Delta$ defined as above, put
\begin{equation}\label{eq:Rankineta}
\mathcal R_{w_0,T}(D_0,\sigma)=D_0^{-\sigma}
\sum_{\substack{w_0<r\le T\\r\ \mathrm{prime}}}
\frac{g(r)}{1-g(r)}r^{3\sigma}
\prod_{\substack{r<p\le T\\p\ \mathrm{prime}}}
\frac{1+g(p)p^\sigma}{1-g(p)},\qquad g(p)=\frac1{p-1}.
\end{equation}
The sum and product in this expression use \emph{all} medium odd primes,
including any primes dividing $N$. Then
\[
0\le\Delta\le\mathcal R_{w_0,T}(D_0,\sigma)V_0,
\qquad V_0\le B^+\le(1+\mathcal R_{w_0,T}(D_0,\sigma))V_0.
\]
\end{lemma}
\begin{proof}
For any finite prime set, expand its full inclusion--exclusion product.
Partition the terms omitted by $b^+$ according to their first failing
odd prefix, and those omitted by $b^-$ according to their first failing
even prefix. Summing all possible remaining primes below the minimum
prime $r$ of that prefix gives $V_0(r)=\prod_{p\mid P_0,\ p<r}(1-g(p))$.
If $T_n^{(0)}$ is the sum of $g(p_1\cdots p_n)V_0(p_n)$ over first
failing prefixes of length $n$ of the relevant parity, this proves the
finite polynomial identities
\[
B^+=V_0+\sum_{n\ \mathrm{odd}}T_n^{(0)},\qquad
B^-=V_0-\sum_{n\ \mathrm{even}}T_n^{(0)},\qquad
\Delta=\sum_{n\ge1}T_n^{(0)}.
\]
They are the finite version of the identities used in BJS Proposition 13.
They also hold when $g(p)=1_{p\mid n}$: all terms are nonnegative,
giving the pointwise brackets. To see the support statement, if the
last prefix has the checked parity then $dp_n^2<D_0$; otherwise, for
$n\ge2$, the preceding checked prefix gives
$d<p_1\cdots p_{n-1}p_{n-1}^2<D_0$. The remaining case is the
one-prime lower weight, where $p\le T<D_0$.

Every first failure satisfies $dp_n^2\ge D_0$. Drop the preceding
prefix restrictions, apply
$1_{dp_n^2\ge D_0}\le(dp_n^2/D_0)^\sigma$, and group by $r=p_n$.
Since
\[
\frac{g(d)V_0(r)}{V_0}
=\frac{g(r)}{1-g(r)}
\prod_{p\mid d,\ p>r}g(p)
\prod_{p\mid P_0,\ p>r}(1-g(p))^{-1},
\]
summing every subset of primes larger than $r$ gives exactly the
positive expression in \eqref{eq:Rankineta} for the actual prime set.
There is no factor two: the odd and even sums partition the prefix
lengths in $\Delta$. Restoring deleted primes adds nonnegative summands
and factors $(1+g(p)p^\sigma)/(1-g(p))\ge1$, so the complete-prime
expression dominates it. Finally $B^+-V_0\le\Delta$.
\end{proof}

\begin{lemma}[Separate finite gaps]\label{lem:parity}
Write $a_0=(B^+-V_0)/V_0$ and $b_0=(V_0-B^-)/V_0$.
For each medium prime $r$, let
\[
H_r(t)=\prod_{r<p\le T}(1+g(p)p^\sigma t),\qquad
e_j(r)=[t^j]H_r(t),\qquad
W_r=\prod_{r<p\le T}(1-g(p))^{-1}.
\]
Both products use all medium primes. At $D_0=T^u$, a prefix with
$k$ primes larger than its minimum $r$ satisfies
$dp_{\min}^2\le T^kr^3$. Retain only degrees for which
$T^kr^3\ge T^u$. If $u>3$, also retain only $k\ge\lfloor u-2\rfloor$.
Define $E_r$ and $O_r$ as the sums of $e_k(r)$ over the retained even
and odd degrees, respectively. Then
\[
a_0\le R_+:=D_0^{-\sigma}\sum_{w_0<r\le T}
\frac{g(r)}{1-g(r)}r^{3\sigma}W_rE_r,\qquad
b_0\le R_-:=D_0^{-\sigma}\sum_{w_0<r\le T}
\frac{g(r)}{1-g(r)}r^{3\sigma}W_rO_r.
\]
The bounds remain valid upon deleting primes dividing $N$.
All terms and elementary coefficients are nonnegative.
\end{lemma}
\begin{proof}
Apply the first-failure proof above separately to odd and even prefix
lengths. An odd prefix has an even number of primes larger than $r$,
and conversely for an even prefix. The degree restriction from
$T^kr^3<D_0$ removes only impossible failures. When $u>3$, a prefix
of length $n$ with $n+2\le u$ cannot fail: for $n=1$ its product is
at most $T^3<T^u$, and for $n\ge2$ distinct primes make the comparison
strict even at equality $n+2=u$. Thus $k=n-1\ge\lfloor u-2\rfloor$.
Each retained polynomial has positive coefficients; restoring primes
increases it and $W_r$, and adds positive minimum-prime terms.
For rational $u=A/B$, the degree test is exactly
$r^{3B}T^{kB}\ge T^A$, so no approximate cutoff is required.

The finite certificate maintains low elementary coefficients and the
two tails beyond a fixed degree $M\ge\lceil u\rceil$.
On adding a larger prime with $h=g(p)p^\sigma$, the low coefficients
update by $e_j'=e_j+he_{j-1}$ in descending order.
If $M$ is odd, the tails update by
$E'=E+hO$ and $O'=O+h(E+e_{M-1})$; if $M$ is even they update by
$E'=E+h(O+e_{M-1})$ and $O'=O+hE$, using old values throughout.
Adding the required low coefficients to these tails evaluates
$E_r,O_r$ using positive arithmetic only.
\end{proof}

\begin{lemma}[Local exponents for separate degrees]\label{lem:degreewise}
Let $\Sigma\subset[0,\infty)$ be finite and nonempty, let $D_0=T^u$
with $u>3$, and let $M\ge\lceil u\rceil$ be an integer.
For each medium prime $r$, partition the retained even and odd degrees
of Lemma \ref{lem:parity} into singleton degrees below $M$ and their
respective residual tails at or above $M$.
For such a group $J$ and $\sigma\in\Sigma$, put
\[
q_{r,\sigma}(J)=D_0^{-\sigma}\frac{g(r)}{1-g(r)}
r^{3\sigma}W_r
\sum_{k\in J}[x^k]\prod_{r<p\le T}(1+g(p)p^\sigma x).
\]
Then
\[
a_0\le\widehat R_+
 :=\sum_r\sum_{J\ {\rm even}}\min_{\sigma\in\Sigma}q_{r,\sigma}(J),
\qquad
b_0\le\widehat R_-
 :=\sum_r\sum_{J\ {\rm odd}}\min_{\sigma\in\Sigma}q_{r,\sigma}(J).
\]
The expressions use complete prime sets and remain valid when primes
dividing $N$ are deleted. Their evaluation uses the positive recurrence
of Lemma \ref{lem:parity} separately for each exponent.
\end{lemma}
\begin{proof}
Fix $r$ and a degree group $J$. Its contribution to the normalized
first-failure sum is independent of $\sigma$. For every $\sigma\ge0$,
the inequality
$1_{dr^2\ge D_0}\le(dr^2/D_0)^\sigma$ bounds that contribution by
$q_{r,\sigma}(J)$; the exponent zero is also permitted.
Every member of the finite family is thus an upper bound for the same
nonnegative contribution, so their minimum is an upper bound.
Sum over the disjoint groups and minimum primes.
Since $M\ge\lceil u\rceil$, all degrees in the residual tails satisfy
the necessary degree conditions; no possible failure is discarded.
For a deleted prime set the complete-prime expression dominates its
corresponding contribution for each exponent. Taking a minimum
preserves this domination. Neither the sieve weights nor their support
have changed; only their finite mass estimates have been tightened.
\end{proof}

\begin{lemma}[Nonnegative quadratic product bounds]\label{lem:quadratic}
Fix a minimum prime $r$ and one of the degree groups $J$ in
Lemma \ref{lem:degreewise}. For $\sigma>0$, use the complete larger
prime set and put
\[
m_j=\sum_{\substack{S\subset\{p:r<p\le T\}\\ |S|\in J}}
 \left(\prod_{p\in S}g(p)\right)
 \left(\frac{r^3\prod_{p\in S}p}{D_0}\right)^{j\sigma},
 \qquad j=0,1,2.
\]
For every $0\le b<1$, the normalized first-failure contribution
of this group is at most
\[
\frac{g(r)}{1-g(r)}W_r
 \frac{m_2-2bm_1+b^2m_0}{(1-b)^2}.
\]
If $m_2<m_1<m_0$, $m_0-2m_1+m_2>0$, and
$b_*=(m_1-m_2)/(m_0-m_1)\in(0,1)$, one may use the bound
\[
\frac{g(r)}{1-g(r)}W_r
 \frac{m_0m_2-m_1^2}{m_0-2m_1+m_2}.
\]
Take a finite minimum with the ordinary Rankin bounds and sum
over the same degree groups and minimum primes. Denote the resulting
gap bounds by $\widehat R^{(2)}_+$ and $\widehat R^{(2)}_-$.
They remain valid when primes dividing $N$ are deleted.
\end{lemma}
\begin{proof}
For $X\ge0$ and $0\le b<1$,
$1_{X\ge1}\le (X-b)^2/(1-b)^2$: the right side is nonnegative
everywhere and at least one on $X\ge1$.
Use $X=(r^3\prod_{p\in S}p/D_0)^\sigma$.
Dropping preceding-prefix restrictions adds only nonnegative
polynomial values. For a deleted prime set, restoring primes
increases both the subset sum of these nonnegative values and $W_r$,
so the displayed complete-prime bound dominates the actual contribution
for every fixed $b$. Only after this domination do we optimize $b$;
no monotonicity of a moment quotient under deletion is assumed.
Differentiating the rational quadratic gives $b_*$ and the second
displayed expression. Finite weighted Cauchy--Schwarz gives
$m_0m_2-m_1^2\ge0$.
The moments at exponents $0,\sigma,2\sigma$ are evaluated by the
positive coefficient and parity-tail recurrences above.
Interval subtraction and division enclose the quadratic expression;
the optimizer is used only when all required signs are certified
strictly. An unresolved comparison is skipped. Enclosing the maximum
of the variance expression and zero is legitimate by its proved
nonnegativity, even when a rounded enclosure of exact zero has a
slightly negative lower endpoint. The weights and support are unchanged.
\end{proof}

\begin{lemma}[A preceding-pass product band]\label{lem:prefixband}
Let $D_0=T^u$ with $u>3$. Every first-failing prefix $d$,
with minimum prime $r$, satisfies
\[
D_0/r^2\le d<D_0.
\]
Consequently $\delta=\log(dr^2/D_0)$ belongs to $[0,R)$,
where $R=2\log r$. For $\sigma,c,b\ge0$, the function
\[
P_{\sigma,c,b}(\delta)
 =\exp(\sigma\delta)\{1+c\delta+b\delta(R-\delta)\}^2
\]
is nonnegative on the real line and at least one on $[0,R]$.
For a retained degree group $J$, take the complete larger prime set
and put
\[
m_j=\sum_{\substack{S\subset\{p:r<p\le T\}\\ |S|\in J}}
 \left(\prod_{p\in S}g(p)\right)e^{\sigma\delta_S}\delta_S^j,\quad
\delta_S=\log(r^3\prod_{p\in S}p/D_0),\quad 0\le j\le4.
\]
Write
\[
v_0=m_1,\quad v_1=Rm_1-m_2,\quad
A=m_2,\quad B=Rm_2-m_3,\quad
C=R^2m_2-2Rm_3+m_4.
\]
The group's normalized first-failure contribution is bounded by
$g(r)W_r/(1-g(r))$ times any of the following certified candidates:
$m_0$; $m_0-v_0^2/A$ if $v_0<0$, $A>0$;
$m_0-v_1^2/C$ if $v_1<0$, $C>0$; and
$m_0+v_0c_*+v_1b_*$ if
\[
AC-B^2>0,\qquad
c_*=\frac{Bv_1-Cv_0}{AC-B^2}\ge0,\qquad
b_*=\frac{Bv_0-Av_1}{AC-B^2}\ge0.
\]
Take finite minima with ordinary Rankin candidates, then sum over
the degree groups and $r$. Denote the resulting gap bounds by
$\widehat R^{\rm band}_\pm$. They remain valid upon deleting primes
dividing $N$; neither the weights nor their support changes.
\end{lemma}
\begin{proof}
A one-prime prefix cannot fail since $T^3<D_0$.
A two-prime prefix has product less than $T^2<D_0$.
For a longer first failure of length $n$, the preceding checked
prefix of length $n-2$ passed. Writing its last prime as $s$ and
the last two primes as $q>r$, we have $s>q>r$ and
$d=d_{n-2}qr<d_{n-2}s^2<D_0$.
Failure gives $dr^2\ge D_0$, proving the band.
On the band both $\delta$ and $\delta(R-\delta)$ are nonnegative;
this proves the majorant. Restoring primes increases the expectation
of this nonnegative function and $W_r$, before any optimization.
This proves domination for every fixed $\sigma,c,b$, including for
deleted prime sets.

The complete expectation of the square equals
$m_0+2cv_0+2bv_1+c^2A+2cbB+b^2C$.
The two boundary candidates and the interior candidate follow by
minimizing this quadratic over $c,b\ge0$ when their stated signs hold.
The matrix with entries $A,B,C$ is a finite weighted Gram matrix,
so all accepted minimum values are nonnegative. Interval maximum
with zero therefore preserves enclosure. All optimizer conditions
are checked strictly; unresolved candidates are omitted.

Raw logarithmic moments are computed positively. Put
\[
E_{k,j}=\sum_{|S|=k}\prod_{p\in S}g(p)p^\sigma
\frac{(\sum_{p\in S}\log p)^j}{j!}.
\]
Adding a prime of logarithm $\ell$
updates $E_{k,j}$ by adding
$g(p)p^\sigma\sum_{a=0}^j\ell^{j-a}E_{k-1,a}/(j-a)!$.
The two parity tails use the same positive convolution on their
opposite tail and the boundary degree, as in Lemma \ref{lem:parity}.
After these positive recurrences, finite binomial shifts give the
$m_j$ with Arb inclusion. The signed shifts and Gram formulas are
not evaluated by floating subtraction or unproved truncation.
Native polynomial multiplication truncated after order four evaluates
these exact moment coefficients: higher orders cannot enter them.
Any fixed logarithm base greater than one gives the same majorant
family after rescaling $\sigma,c,b$ by positive constants.
\end{proof}

\begin{lemma}[Positive prefix blocks]\label{lem:prefixblocks}
Use the finite Rosser weights of Lemma \ref{lem:rankin}, with
$D_0=T^u$ and $u>3$. Let $r_1<r_2<\cdots$ be the complete medium
primes, and let $K$ be the largest integer with
$r_1\cdots r_K<D_0$. Choose $0<h<\log r_1$, put
$J=\lceil\log D_0/h\rceil$, and label a chosen prefix by
$B=\sum_{p\mid d}\lfloor\log p/h\rfloor$.
Partition the complete medium primes into descending blocks with
common label $b=\lfloor\log p/h\rfloor$. In a block set
\[
c=\left\lceil\frac{\log D_0-3\log p_{\min}}h\right\rceil,\qquad
f=\max\left\{0,
\left\lceil\frac{\log D_0-3\log p_{\max}}h\right\rceil-K+1\right\}.
\]
For each of the upper and lower weights, keep two nonnegative
polynomials $X_0,X_1$ for the parity of possible passing prefixes,
truncated to degrees $B<J$, initially $(1,0)$.
Let $A$ append label $b$, flip parity, and, when the appended prefix
has the checked parity, retain only source labels $B<c$.
Always truncate resulting labels to $B<J$.
Let $\mathcal F$ sum source coefficients with $B\ge f$ in the
parity that is checked upon appending one prime.
For $t_p=1/(p-2)$ form the polynomials in the operator variable $a$
\[
D(a)=\prod_{p\ {\rm in\ block}}(1+t_p+t_pa),\qquad
C(a)=\frac{D(a)-1}{1+a}.
\]
Their coefficients are nonnegative. For each block add
$\mathcal F\sum_{j=0}^K[a^j]C(a)A^jX$ to the corresponding gap
bound, and replace $X$ by
$\sum_{j=0}^K[a^j]D(a)A^jX$.
The resulting two totals bound $(B^+-V_0)/V_0$ and
$(V_0-B^-)/V_0$ for every deletion of medium primes.
The rounding changes only the mass bound, not the defining sieve weights.
\end{lemma}
\begin{proof}
Every accepted prefix has product below $D_0$, by the support proof
in Lemma \ref{lem:rankin}; every first-failing prefix does also,
by Lemma \ref{lem:prefixband}. Both therefore have at most $K$
chosen primes. For a source prefix of degree $k$ and true logarithm
$v$, floor labels give $hB\le v\le h(B+k)$.
A checked pass requires $v+3\log p<\log D_0$, so $B<c$.
A first failure has $k\le K-1$ and
$v+3\log p\ge\log D_0$, so $B\ge f$.
Thus these are necessary possible-pass and possible-failure
conditions. A rounding ambiguity may contribute to both; this only
increases the nonnegative bounds. Truncating labels at $J$ removes
no actual accepted prefix.

Before a prime $p$, multiply the accepted-prefix generating
polynomials by the factor
$W=\prod_{q>p}(1-g(q))^{-1}$ over previously processed primes.
For a common relaxed append operator $A$, processing $p$ gives
\[
X'=((1+t_p)I+t_pA)X,\qquad
G'=G+t_p\mathcal FX,
\]
where $G$ is the accumulated normalized gap.
The actual passing prefixes and first failures are bounded by this
positive evolution. The augmented update matrices commute because
$A$ and $\mathcal F$ are common throughout the block. Multiplying
them gives $X'=D(A)X$ and
$G'=G+\mathcal FC(A)X$.
Equivalently, as a prime with value $t$ is added,
\[
D_{\rm new}=(1+t+ta)D,\qquad
C_{\rm new}=(1+t+ta)C+t.
\]
Starting from $D=1,C=0$, these recurrences use positive arithmetic.
With $e_m$ the elementary symmetric sums of the block's $t_p$,
\[
[a^j]D=\sum_{m\ge j}\binom mj e_m,\qquad
[a^j]C=\sum_{m\ge j+1}\binom{m-1}j e_m.
\]
An $A^j$ term selects $j$ primes within this block. Actual
passing prefixes have total degree at most $K$, and actual first
failures have at most $K-1$ preceding selected primes. Hence
omitting terms $j>K$ cannot remove an actual contribution.
Artificial paths of greater total degree may remain; their positive
mass only enlarges the bound.

For a deleted prime set use the same complete-set $K$, blocks,
and cutoffs. They still contain every actual possible pass or failure.
Deleting primes decreases every coefficient of $D$ and $C$,
as their displayed positive formulas show. The positive operators
therefore give complete-set domination. This proves the uniform
claim; no cancellation of prefix masses is used.
\end{proof}

\begin{lemma}[Sharper finite support]\label{lem:sharpsupport}
For the complete medium prime set, let $p_0$ be its least member and put
\[
\kappa=\max\left\{\frac1{p_0^2},
 \max_{\substack{w_0<r<s\le T\\r,s\ {\rm prime}}}\frac r{s^2}\right\},
\]
with the inner maximum omitted for a singleton set.
If $D_0=T^u$, $u>3$, every nonzero coefficient of $b^+$ or $b^-$
has support $d<\kappa D_0$. This remains true for every deleted
prime set. Consequently the composite weights below have support
less than $\kappa Q_*D_0D$.
Their finite gap estimates still use the original level $D_0$.
\end{lemma}
\begin{proof}
If the final prefix is checked, $dr^2<D_0$, hence
$d<D_0/r^2\le\kappa D_0$.
If it is unchecked and has length at least two, the preceding
checked prefix has last prime $s>r$ and passed; thus
$d=d_{\rm prev}r<D_0r/s^2\le\kappa D_0$.
The remaining one-prime lower coefficient satisfies
$d\le T<\kappa T^u$, since $\kappa\ge1/T^2$ and $u>3$.
The empty coefficient also satisfies the bound.
The complete-set maximum dominates every pair in a deleted set.
For a fixed $s$ its largest $r$ is the preceding complete prime,
so a complete adjacent-prime scan computes the maximum.
The convolution support follows by multiplying the exact, medium,
and large support bounds. No coefficient or gap mass has changed.
\end{proof}
For the present cutoffs the exact maxima are $11/169$ at $w_0=5$
and $5/49$ at $w_0=3$. The support-only certificate verifies these
maxima and 52488 finite coefficient/deletion cases. The current
endpoint charges below use $H_0=\kappa D_0$. Their finite gap
estimates still use the original defining level $D_0$.

Take the usual Rosser coefficients $\lambda_1^\pm$ on $P_1$ at
level $D$. They equal either $\mu(d)$ or zero, have support $d<D$,
and their divisor sums bracket $1_{(n,P_1)=1}$; see the construction
used in BJS Proposition 13 and \cite[Theorem 9.3]{Nathanson}. In particular,
the upper divisor sum is nonnegative. A precise definition takes
$d=p_1\cdots p_r$, $p_1>\cdots>p_r$, and requires
$p_1\cdots p_m p_m^2<D$ at every odd $m\le r$ for the upper
weights and every even $m\le r$ for the lower weights.
Their main masses $G_1^\pm$ obey the BJS bounds with $Q=1$, since
either Lemma \ref{lem:uniform} or Lemma \ref{lem:ordered} verifies the
product premise at every endpoint. Write $\mathcal K_T$ for the
uniform product constant supplied by the chosen lemma.

\begin{lemma}[Composite presieving]\label{lem:composite}
Suppose $0<\epsilon\le1/74$, $\mathcal K_T<1+\epsilon$, and use the BJS
constants $C_1,C_2$.
Use either the Brun medium weights, with $H_0=T^k$ and $\eta$ as in
\eqref{eq:Bruneta}, or the finite Rosser medium weights, with
$H_0=D_0$ (or $H_0=\kappa D_0$ under Lemma \ref{lem:sharpsupport})
and $\eta\ge\mathcal R_{w_0,T}(D_0,\sigma)$.
For a nonnegative finite sequence with local density $g$, write
$s=\log D/\log z$ and
$R=\sum_{d\mid P_N(z),\ d<Q_*H_0D}|r(d)|$.
If $D\ge z^2$ and $f(s)-\epsilon C_2e^2h(s)>0$, then the lower
bound in \eqref{eq:sieve} holds. If $D\ge z$, its upper bound holds.
\end{lemma}
\begin{proof}
First convolve the medium/large-prime coefficients to form
\[
\widetilde\lambda^+=b^+*\lambda_1^+,\qquad
\widetilde\lambda^-=b^+*\lambda_1^--(b^+-b^-)*\lambda_1^+.
\]
For pointwise upper and lower brackets $a^-\le I_0\le a^+$
and $c^-\le I_1\le c^+$ with $a^+,c^+\ge0$, one has
\[
a^-c^++a^+c^--a^+c^+\le I_0I_1\le a^+c^+.
\]
For the lower inequality, its left side is at most
$I_0c^++a^+I_1-a^+c^+=I_0I_1-(a^+-I_0)(c^+-I_1)$.
Thus these convolutions are legitimate bounds for the medium/large
parts. Multiply the pointwise inequalities by
$1_{(n,Q_0)=1}\ge0$, equivalently convolve both coefficients with
$\mu_E(d)=\mu(d)1_{d\mid Q_0}$, to obtain the full weights $\lambda^\pm$.
Every divisor has a unique medium/large factorization. Write
$b^\pm(d_0)=\mu(d_0)I^\pm$ and
$\lambda_1^\pm(d_1)=\mu(d_1)J^\pm$, where all four indicators
belong to $\{0,1\}$. The normalized lower coefficient is
$I^+J^--(I^+-I^-)J^+$. If $J^+=0$, it is $I^+J^-$;
if $J^+=1$, it is $I^-$ when $J^-=1$ and $I^--I^+$ otherwise.
Thus it belongs to $\{-1,0,1\}$, as does the normalized upper
coefficient. This argument applies to both medium constructions.
The exact factorization is also unique and $|\mu_E|\le1$, so the full
coefficients have the same bound. Their support is below $Q_*H_0D$.
Consequently no new divisor multiplicity is charged
to $R$.
Writing $V_E=\prod_{p\mid P_E}(1-g(p))$, the upper main mass is
$V_EB^+G_1^+$ and is at most
$(1+\eta)V(z)(F(s)+\epsilon C_1e^2h(s))$.
The lower main mass is $V_E(B^+G_1^--\Delta G_1^+)$ and is at least
$V(z)(f(s)-\epsilon C_2e^2h(s)
-\eta(F(s)+\epsilon C_1e^2h(s)))$.
Here the positivity assumption permits using $B^+\ge V_0$ in
the first term. Expanding the divisor sums against the sequence
and using the coefficient bounds proves the remainder claim.
\end{proof}

The resulting inequalities are
\begin{align}\label{eq:sieve}
S(A,z)&\ge|A|V(z)\{f(s)-\epsilon C_2 e^2h(s)
-\eta(F(s)+\epsilon C_1e^2h(s))\}-R,
&&D\ge z^2,\notag\\
S(A,z)&\le|A|V(z)(1+\eta)\{F(s)+\epsilon C_1 e^2h(s)\}+R,
&&D\ge z,
\end{align}
where the lower bound includes the stated positivity condition, and
$V(z)=\prod_{p\mid P(z)}(1-g(p))$.

When $a_0\le\eta_+$ and $b_0\le\eta_-$ are supplied by Lemmas
\ref{lem:parity}, \ref{lem:degreewise}, \ref{lem:quadratic},
\ref{lem:prefixband}, or \ref{lem:prefixblocks}, put
$l(s)=f(s)-\epsilon C_2e^2h(s)$ and
$u(s)=F(s)+\epsilon C_1e^2h(s)$. Since $u(s)\ge l(s)$ and $u(s)>0$,
the same composite weights give the sharper bounds
\begin{align}\label{eq:sieveparity}
S(A,z)&\ge |A|V(z)\{(1+\eta_+)l(s)-(\eta_++\eta_-)u(s)\}-R,
&&D\ge z^2,\notag\\
S(A,z)&\le |A|V(z)(1+\eta_+)u(s)+R,
&&D\ge z.
\end{align}
Indeed their normalized lower mass is at least
$(1+a_0)l-(a_0+b_0)u=l-a_0(u-l)-b_0u$,
which decreases in both $a_0,b_0$; the upper mass uses only $a_0$.
The earlier coarse form remains available without separate gap inputs.
For the finite Rosser construction these refined inequalities use the
separate inputs of Lemmas \ref{lem:parity}, \ref{lem:degreewise},
\ref{lem:quadratic}, \ref{lem:prefixband}, or \ref{lem:prefixblocks}
directly; the coarse constant can always be chosen as
$\mathcal R_{w_0,T}(D_0,\sigma)$ and is not charged in this refined form.

\section{Analytic inputs and elementary envelopes}
In the ranges needed here,
\[
F(s)=\frac{2e^\gamma}{s}\ (1\le s\le3),\quad
f(s)=\frac{2e^\gamma\log(s-1)}s\ (2\le s\le4),
\]
\[
h(s)=e^{-2}\ (1\le s\le2),\quad e^{-s}\ (2\le s\le3),
\quad 3e^{-s}/s\ (s\ge3).
\]
For $3\le s\le4$, we also retain the exact expression from
\cite{BJS} (the explicit formulas for $F,f$), equivalently obtained
from $(sF(s))'=f(s-1)$:
\begin{equation}\label{eq:F34}
F(s)=\frac{2e^\gamma}{s}
\left(1+\int_3^s\frac{\log(x-2)}{x-1}\,dx\right).
\end{equation}
The integral is evaluated by Arb/Acb analytic integration.
Table 1 of \cite{BJS} permits $(C_1,C_2)=(106,108)$ when
$\epsilon\le10^{-4}$, $(109,110)$ when $\epsilon\le1/2000$,
$(113,114)$ when $\epsilon\le1/1000$,
and $(116,117)$ when $\epsilon\le1/700$.

The prime sequence is
$A=\{N-p:p<N,\ p\text{ prime},\ p\nmid N\}$,
and $g(p)=1/(p-1)$ on $\PP$.
The exact product identity and \eqref{eq:mertens} imply
\begin{equation}\label{eq:Vbounds}
\frac{U_N}{\log x}
\left(1-\delta(x)-\frac{2L}{x}\right)
<V(x)
<\frac{U_N}{\log x}(1+v),
\end{equation}
for the lower bound whenever $x\ge B^{1/3}$, and for the upper
bound at $x=N^{1/8},N^{1/3}$ when $L\ge8500$, with $v=10^{-10}$.
Indeed, the missing factors for $p\mid N$, $p\ge x$, have loss
at most $\om(N)/(x-1)\le L/(\log2\,(x-1))<2L/x$.
The tail $C(x)/C_\infty$ is at most $\exp(1/(x-2))$.
At both large endpoints its effect, the missing-factor loss, and
$\delta(x)$ together are bounded by $Le^{-L/16}<v$.

Schoenfeld's RH bound, for the large $x$ used here, is
$|\pi(x)-\li(x)|\le\sqrt x\log x/(8\pi)$.
For $t=\log x\ge20$ elementary envelopes give
\[
\frac{x}{t}(1+1/t)<\li(x)<\frac{x}{t}(1+1/t+3/t^2).
\]
For completeness, the two inequalities are checked at $t=20$ by
Arb evaluation of $\operatorname{Ei}(20)$; their derivatives differ
from $1/\log x$ by $-2/t^3$ and $(t-9)/t^4$, respectively.
Thus, after subtracting the primes dividing $N$,
\begin{equation}\label{eq:A}
|A|>N/L\quad(N\ge B),\qquad
|A|\le(1+2/L)N/L\quad(L\ge8500).
\end{equation}
The required error comparisons at $B$ follow already from
$L_B=42.832826035\ldots$ and $\om(N)\le L/\log2$.
Likewise, for every $x\ge y=N^{1/3}$ and $L\ge8500$,
\begin{equation}\label{eq:piupper}
\pi(x)\le(1+4/L)x/\log x.
\end{equation}
The error term in the normalized prime count is at most
$L^2e^{-L/6}/(72\pi)$ and is absorbed by $1/L-27/L^2$.

Under GRH, the AP estimate recorded in \cite[Section 2]{BS}, from
Ernvall-Hyt\"onen--Paloj\"arvi, gives
\[
|\pi(N;d,a)-\li(N)/\varphi(d)|\le q_G(N)\sqrt N\log N
\quad(3\le d\le\sqrt N,\ (a,d)=1).
\]
For $L\ge8500$, $q_G(N)<0.17$. To verify this uniformly, omit the
only negative term of the formula in BS and put
\begin{align*}
q_+(w)={}&0.165+12.683/w+254.980/w^2+2607.854/w^3+11605.056/w^4\\
&+e^{-w/4}\{1.314w+0.092\log w+60.883
+8.250\log w/w+939.260/w\}.
\end{align*}
Every positive term decreases for $w\ge8500$; the certificate gives
$q_+(8500)<0.166496$. Set $\beta=7/2$ and
$H=\sqrt N/L^\beta$. Then
\begin{equation}\label{eq:R}
\sum_{\substack{d\le H\\(d,N)=1}}\mu^2(d)|r(d)|
\le0.12\,N/L^{5/2}=:R_A.
\end{equation}
Indeed, $r(1)=0$, and $d=2$ is excluded because $N$ is even. For
$d\ge3$, $\varphi(d)\ge2$, so centering at $\pi(N)/\varphi(d)$
and deleting primes costs at most
\[
\left(0.17+\frac1{16\pi}+\frac{2e^{-L/2}}{\log2}\right)\sqrt N L.
\]
The coefficient is less than $0.195$. Use the unconditional bound
$\sum_{d\le H}\mu^2(d)\le0.608H$ for $H\ge10^9$, recorded in
BS; the certificate verifies $0.608$ times this coefficient is less
than $0.12$. Separately, counting at most $y=N^{1/3}$ prime moduli gives
\begin{equation}\label{eq:Rprime}
\sum_{\substack{z\le q<y\\q\nmid N}}|r(q)|
\le0.2\,N^{5/6}L=:R_{\rm prime}.
\end{equation}
This avoids charging $\log H$ against the entire sum $R_A$.
We also use the unconditional bound
\begin{equation}\label{eq:phimu}
\sum_{d\le H}\frac{\mu^2(d)}{\varphi(d)}\le1.1\log H\quad(H\ge10^9)
\end{equation}
recorded in \cite[Section 2]{BS}.

We will repeatedly use the following consequence of RH.
\begin{lemma}[Weighted reciprocal sums]\label{lem:weighted}
For $L\ge8500$, $z=N^{1/8}$, and all $x\ge z$,
$|\sum_{p<x}1/p-\log\log x-M|\le\xi=L^{-2}$.
For a nonnegative monotone function $w$ on $[a,b]\subseteq[z,\infty)$,
\[
\sum_{a\le p<b}\frac{w(p)}p
\le\int_a^b\frac{w(t)}{t\log t}\,dt+2\xi\max(w(a),w(b)).
\]
\end{lemma}
\begin{proof}
Use \cite[Corollary 1.3]{LN} and
$\delta(z)=3Le^{-L/16}/(64\pi)<L^{-2}$. Stieltjes integration
by parts bounds the difference by
$\xi(w(a)+w(b)+\operatorname{Var}(w))$.
Both endpoint conventions satisfy the same bound by limits.
\end{proof}

\section{The three terms in Chen's inequality}
Set
\[
\alpha=4215/10^6,\quad t=1/2-\alpha,\quad D=N^t,\quad
\epsilon=1/41200,\quad z=N^{1/8},\quad y=N^{1/3},
\]
\[
s=4-8\alpha=3.966280,\quad k=8(1/6-\alpha),\quad
\delta=10^{-8},\quad \ell=\log(1+\delta).
\]
The exact cutoff is $w_0=5$, so $Q_*=15$ and
$V_*^{(0)}=3/8$. Use the finite Rosser medium weights with
$T=5\cdot10^6$, $D_0=T^{49/8}$,
$\eta_+=1/2768$, and $\eta_-=1/2784$.
Use Lemmas \ref{lem:degreewise} and \ref{lem:prefixband} with $M=12$ and
\[
\Sigma_G=\{0,1/10,7/50,9/50,11/50,7/25,9/25\},\quad
\Gamma_G=\{2/25,3/25,4/25,1/5,6/25\}.
\]
The ordered uniform product certificate gives
$\widehat K_T-1<2.426054\cdot10^{-5}<\epsilon$. Positive moment
recurrences and certified shifted Gram combinations over all 348510
medium primes give
$\widehat R^{(\mathrm{band})}_+<0.000361187945<\eta_+$ and
$\widehat R^{(\mathrm{band})}_-<0.000359148152<\eta_-$,
with the degree restrictions of Lemma \ref{lem:parity}.
Lemma \ref{lem:sharpsupport} gives $\kappa=11/169$.
For $L\ge31000$,
$\alpha L-(7/2)\log L>\log(\kappa Q_*)+(49/8)\log T$,
increasing with $L$. The right side is less than $94.454$;
the certified support slack at the endpoint exceeds $0.015045$.
Therefore $\kappa Q_*D_0D\le H$. Also $D\ge z^2$, $D/y\ge z$,
$D\ge y$, $y\le H$, and $H\ge10^9$.
Every sieve in this section uses \eqref{eq:sieveparity}; its
product premise holds uniformly at all intermediate endpoints.

Nathanson's combinatorial Chen inequality, reproduced in the
preliminaries to sieving in \cite{BS}, is
\begin{equation}\label{eq:Chen}
\pi_2(N)\ge S(A,z)-\frac12\sum_{z\le q<y}S(A_q,z)
-\frac12 S(\mathcal B,y)-2N^{7/8}-N^{1/3}-1.
\end{equation}
Here
$\mathcal B=\{N-p_1p_2p_3:z\le p_1<y\le p_2\le p_3,
p_1p_2p_3<N,(p_1p_2p_3,N)=1\}$.
The extra subtraction of one safely excludes $N-p=1$ under either
convention for almost primes.

\subsection{Lower sieve}
From \eqref{eq:sieveparity}, \eqref{eq:A}, \eqref{eq:Vbounds}, and
\eqref{eq:R}, in units $\Th$,
\begin{equation}\label{eq:low}
\begin{aligned}
S(A,z)/\Th&\ge M_0-\frac{0.12}{U_*\sqrt L},\\
M_0&=8(1-v)\bigl\{(1+\eta_+)(f(s)-108\epsilon e^2h(s))\\
&\qquad -(\eta_++\eta_-)(F(s)+106\epsilon e^2h(s))\bigr\}.
\end{aligned}
\end{equation}

\subsection{The prime-factor sum}
Apply the upper sieve to $A_q$ with $D_q=D/q$.
Primes $q\mid N$ contribute nothing and are omitted in all remainder
sums below; main-term upper bounds may include them.
Then $s_q=8(t-\log q/L)\in[k,3]$, so $F(s_q)=2e^\gamma/s_q$.
We retain the correction $e^2h(s_q)$ rather than replacing it by one.
Put $F_*=2e^\gamma/k+106\epsilon$ and $J=\log(8/3)$.
Since both $F(s_q)$ and $h(s_q)$ increase with $q$, their sum is a
nonnegative monotone weight, bounded by $F_*$.
Lemma \ref{lem:weighted} and
$1/(q-1)\le (1-z^{-1})^{-1}/q$ give
\[
\sum_{z\le q<y}\frac{F(s_q)+106\epsilon e^2h(s_q)}{q-1}
\le\frac{I+106\epsilon J_h(t)+2F_*/L^2}{1-z^{-1}},
\]
where direct integration yields
\[
I=\int_{1/8}^{1/3}\frac{F(8(t-\beta))}{\beta}\,d\beta
=\frac{e^\gamma}{4t}\log\left(6\frac{3-8\alpha}{3-18\alpha}\right).
\]
Writing $\beta_0=t-1/4$, which lies between $1/8$ and $1/3$,
the two pieces of $h$ give
\begin{equation}\label{eq:Jh}
J_h(t)=\int_{1/8}^{\beta_0}\frac{e^{2-8(t-\beta)}}{\beta}\,d\beta
+\log\frac{1/3}{\beta_0}
=\int_{1/8}^{1/3}\frac{e^2h(8(t-\beta))}{\beta}\,d\beta.
\end{equation}
Thus $0<J_h(t)<J$. At the fixed parameter value, analytic interval
integration gives $J_h(t)=0.715653586850489\ldots$.
In the error sum, $r_q(d)=r(qd)-r(q)/\varphi(d)$.
The map $(q,d)\mapsto qd$, for $d\mid P(z)$ and $q\ge z$,
is injective. Thus the total sieve remainder is at most
$R_A+0.55L R_{\rm prime}$ by \eqref{eq:phimu}. The substitution
$|A_q|=|A|/(q-1)+r(q)$ costs at most
$(1+\eta_+)V(z)F_*R_{\rm prime}$. It follows that
\begin{equation}\label{eq:upper}
\frac{1}{2\Th}\sum S(A_q,z)\le M_1+
\frac{0.06}{U_*\sqrt L}
+\frac{0.055L^4e^{-L/6}}{U_*}
+0.8(1+\eta_+)(1+v)F_*L^2e^{-L/6},
\end{equation}
\[
M_1=\frac{4(1+\eta_+)(1+2/L)(1+v)}{1-e^{-L/8}}
\left(I+106\epsilon J_h(t)+\frac{2F_*}{L^2}\right).
\]

\subsection{The switched triple-prime mass}
Partition $p_1$ into the intervals $[w_j,(1+\delta)w_j)$, truncated
at $y$, where $w_j=z(1+\delta)^j$. Enlarge the product restriction
to $w_jp_2p_3<N$, and drop $(p_1,N)=1$. Denote the disjoint enlarged
sets by $\mathcal B_j$, with union $\overline{\mathcal B}$.
Then $\mathcal B\subseteq\overline{\mathcal B}$ and its products
are less than $(1+\delta)N$. Unique sorted prime factorization
ensures that these are legitimate sets, including $p_2=p_3$.
Some enlarged differences $N-p_1p_2p_3$ can be negative. Translate
each finite sequence by a sufficiently large integer multiple of
$P_N(y)$ to make its indices positive. This preserves every divisibility
condition $d\mid P_N(y)$, its cardinality, and all sieve remainders,
so the arithmetic-function form of the cited sieve is still usable.

For $u=p_1$, put $w=\sqrt{(1+\delta)N/u}$ and
\[
H(u)=\int_y^{\sqrt{N/u}}
\frac{dv}{v\log v\log(N/uv)}
=\frac{\log(2-3\beta)}{L(1-\beta)},\quad \beta=\log u/L.
\]
By \eqref{eq:piupper}, the mass is bounded by
$(1+\delta)(1+4/L)N$ times
$\sum_{z\le p_1<y}p_1^{-1}
\sum_{y\le p_2<w}(p_2\log(N/p_1p_2))^{-1}$.
The inner weight is increasing and at most $4/L$.
Lemma \ref{lem:weighted} costs $8/L^3$; extending its integral
from $\sqrt{N/u}$ to $w$ costs at most $6\ell/L^2$.
The function $H(u)$ is positive decreasing, with $H(z)<1/L$,
so the outer weighted sum costs at most $2/L^3$.
Consequently
\begin{equation}\label{eq:mass}
|\overline{\mathcal B}|\le\frac{(1+\delta)(1+4/L)N}{L}
\left\{\bar c+\frac2{L^2}
+\left(J+\frac2{L^2}\right)
\left(\frac{6\ell}{L}+\frac8{L^2}\right)\right\},
\end{equation}
\[
\bar c=\int_{1/8}^{1/3}\frac{\log(2-3\beta)}{\beta(1-\beta)}\,d\beta
=0.3630837292483046\ldots<0.363084.
\]
The integral is certified by Arb/Acb analytic integration, with the
logarithm's analyticity flag enabled. This direct argument uses a
prime-count \emph{upper} envelope and does not use the negative RH
error printed in the triple-mass estimate of the old GRH preprint.

Apply the upper sieve at $y$ to every $\mathcal B_j$ with level $D$.
Its parameter $3t$ lies in $[1,2]$. Therefore, using \eqref{eq:mass},
the normalized half-main-term is at most
\begin{equation}\label{eq:switch}
M_2=\frac{(1+\eta_+)(1+v)}{2}\left(\frac{2e^\gamma}{t}+318\epsilon\right)
(1+\delta)(1+4/L)
\left\{\bar c+\frac2{L^2}+
\left(J+\frac2{L^2}\right)\left(\frac{6\ell}{L}+\frac8{L^2}\right)\right\}.
\end{equation}

\section{Uniform rectangle errors and the half-line certificate}
The GRH bilinear lemma in \cite{BS} states, for $|a(n)|\le1$,
$X,Y,Z\ge25$, $Z<Y$, $H\ge10^9$, and
\[
\left(\frac{\sqrt Y\log H}{q_G(Y)\log Y\log2}\right)^{2/3}\le Z,
\]
that the total square-free AP error of
$\sum_{n<X}\sum_{Z\le p<Y}a(n)$ is at most
\begin{equation}\label{eq:bilinear}
m(X,Y,H)\frac{XY\log^{5/3}H\log^{1/3}Y}{Y^{1/6}}.
\end{equation}
The principal term here restricts $(np,d)=1$. The explicit coefficient is
\begin{align*}
m={}&2.81q_G(Y)^{1/3}
+\frac{2.809}{q_G(Y)^{2/3}\sqrt Y\log Y}
+\frac{1.721\log^{2/3}H}{q_G(Y)^{1/3}Y^{2/3}\log^{2/3}Y\log\log H}\\
&+5.498(X^{-1/2}+Y^{-1/2})Y^{1/6}
\frac{\log^{1/3}H}{\log^{1/3}Y}
+\frac{19.044H\log^{1/3}H}{XY^{5/6}\log^{1/3}Y}.
\end{align*}
For our rectangles, $X=N/w_j$, $Y=\min(y,(1+\delta)w_j)$,
$Z=w_j$, and $a(n)$ indicates $n=p_2p_3$ with
$y\le p_2\le p_3$ and $(n,N)=1$. Thus $|a(n)|\le1$,
$X\ge N^{2/3}$, $Y\ge z$, $Y<\sqrt N$, and $XY\le(1+\delta)N$.
The AP bound gives $0.16<q_G(Y)<0.64$ in this range; this follows
also directly from the formula in \cite[Section 2]{BS}, since
$\log Y\ge1062.5$, its constant term is $0.165$, and its only
negative term is exponentially small.
Termwise, without any monotonicity assertion for $m$, we have
\begin{align*}
m\le{}&2.81(0.64)^{1/3}
+\frac{2.809e^{-L/16}}{(0.16)^{2/3}(L/8)}
+13e^{-L/12}\\
&+9(e^{-5L/18}+e^{-L/24})+31e^{-13L/48}<3.
\end{align*}
The rectangle condition follows from
$13(1+\delta)^{1/3}e^{-L/12}<1$.
Since both logarithms in \eqref{eq:bilinear} are at most $L/2$,
the error in one rectangle is at most
$\frac34(1+\delta)N^{47/48}L^2$.

Removing $(np_1,d)=1$ from the principal term costs at most
$2(1+\delta)N\om(d)/(z\varphi(d))$: a prime dividing $d$ can
divide either $n$ or $p_1$, and both possibilities are counted.
Using $\om(d)\le\log H/\log2$ and \eqref{eq:phimu}, its sum
is at most $(1+\delta)N^{7/8}L^2$, since
$2\cdot1.1/(4\log2)<1$ and $\log H\le L/2$.
There are $\lceil5L/(24\ell)\rceil<L/\ell$ rectangles.
Thus the normalized half-remainder is bounded by
\begin{equation}\label{eq:Erect}
E_{\rm rect}(L)=\frac{(1+\delta)L^5}{2U_*\ell}
\left(\frac34e^{-L/48}+e^{-L/8}\right).
\end{equation}
All endpoints, including the truncated final rectangle, have now
been included; no floating-point underflow is used to discard errors.

Combining \eqref{eq:Chen}--\eqref{eq:Erect}, set
\begin{align*}
E_{\rm AP}(L)&=\frac{0.18}{U_*\sqrt L}
+\frac{0.055L^4e^{-L/6}}{U_*}
+0.8(1+\eta_+)(1+v)F_*L^2e^{-L/6},\\
E_{\rm fin}(L)&=\frac{L^2}{U_*}
\left(2e^{-L/8}+e^{-2L/3}+e^{-L}\right).
\end{align*}
Then
\[
\pi_2(N)/\Th\ge
M_0-M_1(L)-M_2(L)-E_{\rm AP}(L)-E_{\rm rect}(L)-E_{\rm fin}(L).
\]
Every subtracted term is decreasing for $L\ge31000$.
For the exponential terms use $L^be^{-aL}$ decreasing when $L>b/a$;
the largest such cutoff is $240$. All other factors are positive
decreasing rational functions of $L$ or $(1-e^{-L/8})^{-1}$.
The fixed $\alpha$ makes $M_0$ constant and avoids any moving
parameter issue. At $L=31000$, the 192-bit certificate gives
\begin{center}
\begin{tabular}{lr}\toprule
Term & certified value (rounded for display)\\\midrule
$M_0$ & $7.806859399895$\\
$M_1$ & $6.498312561553$\\
$M_2$ & $1.306394071261$\\
$E_{\rm AP}$ & $0.001379428851$\\
$E_{\rm rect}$ & $<10^{-248}$\\
$E_{\rm fin}$ & $<10^{-1670}$\\\midrule
Surplus & $>0.0007733382290$\\\bottomrule
\end{tabular}
\end{center}
This proves the earlier counting coefficient $0.0006$,
using the uniform product premise and composite weights established above.
For its lower starting point, keep all finite weights and gap inputs fixed,
but take $\alpha=4238/10^6$ and $L_0=30824$.
The same 192-bit checks of every support, parameter and remainder
condition give support slack $>0.002084$ and
\[
M_0-M_1-M_2-E_{\rm AP}-E_{\rm rect}-E_{\rm fin}
>0.0000337669646>10^{-5}.
\]
Here $E_{\rm rect}<10^{-248}$ and $E_{\rm fin}<10^{-1660}$;
both are charged positively. The new $s=3.966096$ and $k>1$ obey
all the same analytic ranges. With this fixed $\alpha$ the support slack
increases and every deducted expense decreases on $L\ge30824$,
by the same monotonicity argument. Thus the earlier bound holds for
$N\ge e^{30824}$, with $\log(30824)=10.336048886\ldots$ and about
$13387$ decimal digits at its starting point.
Both calculations cover the entire half-line, rather than a finite mesh.

\subsection{An enlarged finite grid and a sharper AP budget}\label{sec:densegrh}
The defining weights, $T$, $D_0$, $Q_*$, $M$, $\Sigma_G$, and
$\kappa$ remain as above. Enlarge only the finite majorant choices to
\[
\Gamma_G^{\rm new}=\{a/25:2\le a\le9\}.
\]
Lemma \ref{lem:prefixband} applies to each of these nonnegative
majorants; taking their finite minimum cannot increase either gap.
A fresh complete integer sieve and 192-bit positive-moment recurrence
over all 348510 medium primes gives
\[
\widehat R_+<0.000349096808<1/2864,\qquad
\widehat R_-<0.000347781454<1/2875.
\]
Set $\epsilon=1/41219$, which still exceeds the previously certified
uniform product error, and use these two new gap budgets.
With $\beta=7/2$ and the previous AP charge, the endpoint already
falls to $L_0=30779$.

There is a further saving in the AP charge. For a fixed
$7/2\le\beta<4$, set $H=\sqrt N/L^\beta$ and retain the positive
envelope $q_+(L)$ above, rather than replacing it by $0.17$.
The proof of \eqref{eq:R} gives
\[
R_A\le0.608\,P(L)\frac{N}{L^{\beta-1}},\qquad
P(L)=q_+(L)+\frac1{16\pi}+\frac{2e^{-L/2}}{\log2}.
\]
Its coefficients in the lower sieve and half of the prime-factor
sum total $3/2$. Hence replace only the leading term
$0.18/(U_*\sqrt L)$ in $E_{\rm AP}$ by
\begin{equation}\label{eq:denseap}
\frac{0.912\,P(L)}{U_*L^{\beta-3}}.
\end{equation}
Keep the two other AP charges, the entire rectangle charge, and
every finite loss. The support condition becomes
$\alpha L-\beta\log L>\log(\kappa Q_*D_0)$.
All rectangle estimates above used only $H\le\sqrt N$ and
$H\ge10^9$, and continue to hold.

That preceding threshold had $\log(30747)=10.333547707\ldots$
and about $13354$ decimal digits.
For the preceding assertions $N\ge e^{30747}$ with coefficient $10^{-5}$
and $N\ge e^{31000}$ with coefficient $0.001$, use the following
fixed parameters and strict certified endpoint inequalities:
\begin{center}
\begin{tabular}{lrr}\toprule
 & Lower starting point & Stronger count\\\midrule
$L_0$ & $30747$ & $31000$\\
$\alpha$ & $4255/10^6$ & $4218/10^6$\\
$\beta$ & $88/25$ & $351/100$\\
$s$ & $3.965960$ & $3.966256$\\
Support slack & $>0.0005409294$ & $>0.0046277428$\\
$E_{\rm AP}$ & $<0.0010576384$ & $<0.0011678640$\\
$E_{\rm rect}$ & $<10^{-248}$ & $<10^{-250}$\\
$E_{\rm fin}$ & $<10^{-1659}$ & $<10^{-1673}$\\
Surplus & $>0.0000139212861$ & $>0.0010866414727$\\\bottomrule
\end{tabular}
\end{center}
The certificate checks every preceding analytic range, the actual
$H\ge10^9$, and the support inequality at each new endpoint.
For the half-line, $\alpha-\beta/L>0$ and
$1/2-\beta/L>0$ at the endpoint and thereafter. Thus support slack
and $H$ increase. The function $P(L)$ decreases for
$L\ge8500$, and $\beta>3$; \eqref{eq:denseap} therefore decreases.
The remaining monotonicity arguments are unchanged.
This proves those earlier assertions for all $N$ in their respective half-lines.
No distribution hypothesis beyond GRH has been introduced.
The eight tested values $\beta=3.50,3.51,\ldots,3.57$ are a finite
parameter search, not a proof of optimality.

% BEGIN CURRENT PREFIX GRH
\subsection{All preceding checks and the actual modulus count}\label{sec:prefixgrh}
The following calculation proves the earlier starting point $e^{26820}$
and the preceding stronger coefficient $0.0179$.
Use Lemma \ref{lem:prefixblocks} with $T=5000000$, $w_0=5$,
$Q_*=15$, $u=87/16$, and $h=1/200$.
The complete integer sieve has 348510 medium primes;
the exact degree bound is $K=22$, the label limit is $J=16775$,
and there are 1810 nonempty blocks. The strict positive recurrences give
\[
\widehat R_+<0.000369501096<1/2706,\qquad
\widehat R_-<0.000369035200<1/2709.
\]
Set $\epsilon=1/41219$, $\eta_+=1/2706$, $\eta_-=1/2709$,
$(C_1,C_2)=(106,108)$, and retain $\kappa=11/169$.
The actual support cost is
\[
c=\log(\kappa Q_*)+u\log T=83.849204066\ldots.
\]
All uses of \eqref{eq:sieveparity} accept the new separate gap inputs;
its proof depends only on their domination of the two finite masses.

\begin{lemma}[Counting the actual modulus support]\label{lem:sparsecount}
For the composite coefficients in Lemma \ref{lem:composite}, put
$H_0=\kappa D_0$ and let $m$ be the number of complete exact primes.
For every $\sigma>0$, the number of their nonzero coefficients is at most
\[
2^m D H_0^\sigma
\prod_{w_0<p\le T}(1+p^{-\sigma}).
\]
This bound holds for every deletion of primes dividing $N$ and both signs.
If $Q_*H_0D\le H$, this number is at most $\nu H$, where
\begin{equation}\label{eq:sparsefraction}
\nu=\frac{2^m}{Q_*}H_0^{\sigma-1}
\prod_{w_0<p\le T}(1+p^{-\sigma}).
\end{equation}
\end{lemma}
\begin{proof}
Every contributing divisor has a unique exact, medium and large factorization.
There are at most $2^m$ exact divisors. Every medium factor is square-free,
uses only primes in $(w_0,T]$, and is below $H_0$ by Lemma
\ref{lem:sharpsupport}, for either medium sign. Rankin's elementary inequality
$1_{d_0<H_0}\le(H_0/d_0)^\sigma$ bounds their number by the Euler product
in the statement. A large factor is a positive integer below $D$,
so there are fewer than $D$ choices. Counting all such triples bounds
the coefficient support, even when some coefficients vanish.
Deleting primes decreases the positive Euler product and the available
exact factors, while the complete-set $H_0$ still bounds their support.
Dividing by $Q_*H_0D\le H$ proves the second assertion.
\end{proof}

For $\sigma=4/5$, a fresh 192-bit Euler-product sum over the complete
medium primes gives
\[
\log\prod_{5<p\le T}(1+p^{-\sigma})<10.1391321865,\qquad
\nu<0.000604653842<1/1653.
\]
For the prime-factor sieves $A_q$, $q\ge z>T$ and the large level is $D/q$.
For their aggregate AP remainder, write $d=e d_0d_1$ and absorb $q$
into the large factor: $qd_1<D$. Since $d\mid P(z)$, the prime $q$
is the unique prime factor of $qd$ that is at least $z$. Thus
$(q,d)\mapsto qd$ is injective over the whole sum. All its images
belong to the same enlarged set of exact, medium and large triples,
so the same count bounds the aggregate first AP remainder.
The additional centering term
$r(q)/\varphi(d)$ retains its previous charge.

Set $H=\sqrt N/L^\beta$ with $\beta\ge3$. The pointwise AP bound and
centering argument used for \eqref{eq:R} give, for each actual support,
\[
R_{\rm support}\le\nu P(L)\frac{N}{L^{\beta-1}},\qquad
P(L)=q_+(L)+\frac1{16\pi}+\frac{2e^{-L/2}}{\log2}.
\]
Thus the lower sieve and half of the prime-factor sum charge a total of
$3R_{\rm support}/2$. Replace \eqref{eq:denseap} by
\begin{equation}\label{eq:sparseap}
\frac{3\nu P(L)}{2U_*L^{\beta-3}},
\end{equation}
using the fixed upper budget $\nu_*=1/1653$ in this expense.
Both other AP terms, the prime-modulus centering cost, the full rectangle
charge, and every finite loss remain in the proof. The bilinear lemma
still uses the full maximum modulus $H$, rather than this count.

The strict endpoint certificate gives
\begin{center}
\begin{tabular}{lrr}\toprule
 & Lower starting point & Stronger count\\\midrule
$L_0$ & $26820$ & $31000$\\
$\alpha$ & $4267/1000000$ & $1853/500000$\\
$\beta$ & $3$ & $3$\\
$s$ & $3.965864$ & $3.970352$\\
Support slack & $>0.0010264627$ & $>0.0115684828$\\
$E_{\rm AP}$ & $<0.0002269656$ & $<0.0002268874$\\
Surplus & $>0.000034241255$ & $>0.017948372354$\\
$E_{\rm rect}$ & $<10^{-212}$ & $<10^{-250}$\\
$E_{\rm fin}$ & $<10^{-1446}$ & $<10^{-1673}$\\\bottomrule
\end{tabular}
\end{center}
The respective surpluses exceed $10^{-5}$ and $0.0179$.
All analytic ranges used in the earlier calculation are checked again,
including $H\ge10^9$, $T<z$, $3<s<4$, $k>1$, and actual support
$\alpha L-\beta\log L>c$.
For each fixed parameter pair, support slack and $H$ increase thereafter.
With $\beta=3$, the leading AP expense decreases with $P(L)$;
every other expense decreases as before. The inequalities therefore
hold on the whole half-lines. This uses GRH alone.

The previous prefix-block budget using the full square-free modulus count
gave $N\ge e^{28279}$ and coefficient $0.0116$ at $e^{31000}$.
The original finite grid gave $e^{28334}$.
The new support count accounts for the further improvement.
The searched finite parameters and the four Rankin exponents do not
prove global optimality.
% END CURRENT PREFIX GRH
% BEGIN CURRENT ROUGH INTERVAL
\subsection{Logarithmic support counts and rough large factors}\label{sec:roughinterval}
Two elementary count bounds improve the previous support estimate.
They introduce no distribution hypothesis. Throughout this subsection,
the actual sieve coefficients remain those of Lemma \ref{lem:composite}.

\begin{lemma}[Positive logarithmic-bin support count]\label{lem:binnedcount}
Let $\mathcal M$ be the complete medium-prime set, $H_0=\kappa T^u>1$,
and let $K$ be the largest number of its smallest primes whose product
is below $H_0$. Choose $0<h<\log\min\mathcal M$, put
$b(p)=\lfloor\log p/h\rfloor$ and $J=\lceil\log H_0/h\rceil$.
Partition $\mathcal M$ into blocks $\mathcal M_b$ of common label $b$.
For $\sigma\ge0$, define positive coefficients by
\[
\prod_{p\in\mathcal M_b}(1+p^{-\sigma}a)=\sum_j e_{b,j}a^j,
\qquad
P(v)=\prod_b\left(\sum_{0\le j\le K}e_{b,j}v^{jb}\right)\pmod{v^J}.
\]
Write $P_B=[v^B]P$. The number of square-free products of primes in
$\mathcal M$ that are below $H_0$ is at most
\begin{equation}\label{eq:binnedcount}
A_\sigma=\sum_{B<J}P_B
 \exp\!\left(\sigma\min\{\log H_0,h(B+K)\}\right).
\end{equation}
The same complete-set upper bound holds after arbitrary prime deletion.
For $\sigma=0$, $e_{b,j}=\binom{|\mathcal M_b|}{j}$ and $A_0=\sum_{B<J}P_B$.
\end{lemma}
\begin{proof}
A true product $d<H_0$ contains at most $K$ primes. Its label $B$ satisfies
$hB\le\log d<\log H_0$, hence $B<J$. Its block degrees are at most $K$,
so its positive contribution $d^{-\sigma}$ occurs in $P_B$.
Every intermediate label is at most its final label; truncation modulo
$v^J$ after each block cannot remove that contribution. Also
$\log d\le h(B+K)$, so its contribution to \eqref{eq:binnedcount} is at
least one. Products of excessive total degree that survive the block
truncations are harmless additional positive contributions.
Deleting primes decreases every elementary coefficient $e_{b,j}$,
while the complete-set degree and label bounds remain valid.
\end{proof}

\begin{lemma}[Counting rough integers]\label{lem:integerrough}
Let $V_T=\prod_{p\le T}(1-1/p)$, and let $\eta_+$ be the complete
normalized medium upper-gap budget supplied by Lemma \ref{lem:prefixblocks}
for $g(p)=1/(p-1)$. With the same $D_0$, $H_0$ and exact odd product $Q_*$,
the number $R_T(D)$ of positive integers $n<D$ having no prime factor
at most $T$ satisfies
\begin{equation}\label{eq:integerrough}
R_T(D)\le a_TD+b_T,\qquad
a_T=(1+\eta_+)V_T,\qquad b_T=2Q_*H_0.
\end{equation}
\end{lemma}
\begin{proof}
Use full inclusion--exclusion at all primes at most $w_0$, including $2$,
and the finite Rosser upper coefficients at $w_0<p\le T$.
For the integer sequence the density is $g_0(p)=1/p$.
The normalized positive recurrences in Lemma \ref{lem:prefixblocks}
also hold with $t_p=g(p)/(1-g(p))$: their passing and failing operators
depend only on the prime labels and $D_0$, and every coefficient of
$D(a)$ and $C(a)$ is a nonnegative polynomial in the $t_p$.
Thus decreasing $t_p$ decreases every propagated state and gap bound.
Here $t_{0,p}=1/(p-1)\le1/(p-2)$, so the normalized medium upper gap
at density $g_0$ is at most $\eta_+$. Multiplying by the exact small-prime
factor gives an upper main mass $\sum_d\lambda_d^+/d\le a_T$.

These integer upper weights have $|\lambda_d^+|\le1$ and support
$d<2Q_*H_0$ by Lemma \ref{lem:sharpsupport}; the factor $2$ is required
for the exact even prime. The number of positive multiples of $d$
strictly below $D$ is $\lceil D/d\rceil-1$, differing from $D/d$ by
an amount in $[-1,0)$. Summing the pointwise upper weights therefore gives
\[
R_T(D)\le\sum_d\lambda_d^+(\lceil D/d\rceil-1)
\le D\sum_d\frac{\lambda_d^+}{d}+\#\{d:\lambda_d^+\ne0\}
\le a_TD+b_T.
\]
This is an elementary integer-count bound, not an AP distribution input.
\end{proof}

Put $m=\omega(Q_*)$ for the complete exact odd-prime set and
\[
\nu_\sigma=\frac{2^m A_\sigma}{Q_*H_0}.
\]
Every large-prime factor of a composite coefficient has no prime factor
at most $T$. The coefficient support therefore has size at most
$2^mA_\sigma(a_TD+b_T)$. If $Q_*H_0D\le H$, its size is at most
\begin{equation}\label{eq:roughsupport}
\nu_\sigma(a_T+b_T/D)H.
\end{equation}
This counts all possible exact, medium and large factors; any vanishing
coefficient only reduces the support. Prime deletion preserves the bound.
For the aggregate $A_q$ remainder, write $d=e d_0d_1$ with large level
$D/q$ and absorb $q$ into $d_1$. The resulting large factor
$qd_1<D$ has no prime factor at most $T$. Because $d\mid P(z)$ and
$q\ge z$, $q$ is the unique prime factor of $qd$ that is at least $z$.
The map $(q,d)\mapsto qd$ is therefore injective over the whole sum,
and its complete image has at most $2^mA_\sigma R_T(D)$ elements.
Thus \eqref{eq:roughsupport} bounds the aggregate first AP remainder,
without summing a separate floor charge for each $q$.
The separate prime-modulus centering cost is retained.

For $T=5000000$, $w_0=5$, $u=87/16$ and $\kappa=11/169$,
use the preceding gap budgets and $\sigma=0$, $h=1/100$ for the count.
The complete count uses 348510 medium primes, $K=22$, $J=8115$,
and 992 nonempty blocks. Its positive polynomial calculation yields
\[
A_0<5.488142\cdot10^{30},\qquad
\nu_0<0.000008437984<\nu_*:=1/118511.
\]
The gap computation still uses $h=1/200$ as in Section \ref{sec:prefixgrh};
the two bin widths serve different bounds.
An independently regenerated complete integer prime set contains 348513
primes through $T$ and gives
\[
V_T<0.036397999733,\qquad
a_T<a_*:=36411450583/10^{12}.
\]

\paragraph{Finite analytic interval and the existing tail.}
Let $L_*=26820$, fix $\alpha=2117/500000$, $\beta=209/100$,
and put $D=e^{tL}$, $t=1/2-\alpha$, $H=e^{L/2}/L^\beta$.
Set $L_0=24799$ and $c=\log(Q_*H_0)$.
On the closed interval $L_0\le L\le L_*$, the support condition
$\alpha L-\beta\log L>c$ persists since
$\alpha-\beta/L_0>0$. Also $H\ge10^9$, $H\le\sqrt N$,
$T<z$, $3<s<4$ and $k>1$ hold throughout.
Using \eqref{eq:roughsupport} in the pointwise GRH AP bound gives the
normalized leading expense
\[
\frac{3\nu_*}{2U_*}(a_*+b_Te^{-tL})P(L)L^{3-\beta}.
\]
Because $0<\beta<3$, this expression is not claimed to decrease on a
whole half-line. Instead, use the constant upper bound
\begin{equation}\label{eq:intervalplateau}
C_{\rm AP}=\frac{3\nu_*}{2U_*}
(a_*+b_Te^{-tL_0})P(L_0)L_*^{3-\beta}
\end{equation}
throughout this finite interval. This is valid because $P(L)$ and
$e^{-tL}$ decrease, and $L^{3-\beta}\le L_*^{3-\beta}$.
The positive floor charge is kept; at $L_0$,
$b_Te^{-tL_0}<10^{-5302}$.

Replace only the leading AP term by $C_{\rm AP}$.
Both other AP terms described after \eqref{eq:denseap}, every prime-modulus centering
charge, \eqref{eq:Erect}, and the finite loss are retained.
The bilinear estimate still uses the maximum modulus $H$.
Its proof in Section 6 needs $H\le\sqrt N$ and $H\ge10^9$, both verified
here; it does not use the smaller support count.
The endpoint certificate gives $s=3.966128$ and
\begin{center}
\begin{tabular}{lr}\toprule
Term & Certified value\\\midrule
$M_0$ & $>7.806706554276225$\\
$M_1$ & $<6.498953398927134$\\
$M_2$ & $<1.306496415897830$\\
$E_{\rm AP}$ & $<0.001235093952316$\\
Support slack & $>0.001974440924078$\\
Surplus & $>0.000021645498948$\\
$E_{\rm rect}$ & $<10^{-194}$\\
$E_{\rm fin}$ & $<10^{-1337}$\\\bottomrule
\end{tabular}
\end{center}
For these fixed parameters $M_0$ is constant, while $M_1$, $M_2$,
the other AP terms and all remaining expenses decrease for $L\ge L_0$.
Thus the strict surplus exceeding $10^{-5}$ at $L_0$ holds throughout
$[L_0,L_*]$. Section \ref{sec:prefixgrh} already proves the same
claim for every $L\ge L_*$, with common normalization $U_NN/L^2$.
These inclusive ranges cover the whole half-line $L\ge L_0$ and prove
the preceding starting point $e^{24799}$ using GRH alone.
The preceding stronger $0.0179$ coefficient at $L\ge31000$ retains
its earlier half-line certificate.

The separate-formula check reproduces 12 endpoint quantities and
regenerates the integer product with an independent prime sieve.
The histogram checks enumerate 9216 subsets and verify 17172 exact
rational coefficients at $\sigma=0,1$; the integer checks cover 16384
divisor-floor terms and 24316 roughness patterns.
An additional finite enumeration checks 8608 prime-factor modulus patterns,
including the injective aggregate map and its rough-factor support count.
These finite checks accompany the two general proofs above.
The 135 endpoint candidates are a finite parameter grid, not a global
optimality proof. This analytic-regime join does not connect to the
finite Goldbach verification endpoint $4\cdot10^{18}$.
% END CURRENT ROUGH INTERVAL
% BEGIN CURRENT ACCEPTED INTERVAL
\subsection{The union of the accepted medium supports}\label{sec:acceptedinterval}
The preceding logarithmic count includes medium products that neither
Rosser sign accepts. The following positive count keeps every preceding
acceptance condition. It bounds the union of the two supports once,
without changing the sieve coefficients in Lemma \ref{lem:composite}.

\begin{lemma}[Six-state union count]\label{lem:acceptedcount}
Let $D_0=T^u$, $u>3$, and let $H_0=\kappa D_0>1$ bound both pure-medium
Rosser supports as in Lemma \ref{lem:sharpsupport}. Use the complete
medium-prime set, the degree bound $K$, labels $b(p)$ and cutoff $J$
of Lemma \ref{lem:binnedcount}, with $0<h<\log\min\mathcal M$.
Process the complete blocks in decreasing label order. In block
$\mathcal M_b$, put
\[
c_b=\left\lceil\frac{\log D_0-3\log\min\mathcal M_b}{h}\right\rceil,
\qquad m_b=|\mathcal M_b|.
\]
For a positive polynomial $F=\sum_B F_Bv^B$, define the positive operators
\[
\mathcal L_cF=\sum_{B<c}F_Bv^B,\quad
\mathcal H_cF=\sum_{B\ge c}F_Bv^B,\quad
\mathcal S_bF=v^bF\pmod{v^J}.
\]
Track the vector $X=(U_e,U_o,L_e,L_o,B_e,B_o)^\mathsf T$.
The first two coordinates mean only the upper sign is alive, the next
two mean only the lower sign is alive, and the final two mean both signs
are alive; subscripts record the parity of the number of selected primes.
Start with $B_e=1$ and every other coordinate zero. Within the block,
the append operator is
\begin{equation}\label{eq:acceptedoperator}
\mathcal A_bX=
\begin{pmatrix}
\mathcal S_b(U_o+\mathcal H_{c_b}B_o)\\
\mathcal S_b\mathcal L_{c_b}U_e\\
\mathcal S_b\mathcal L_{c_b}L_o\\
\mathcal S_b(L_e+\mathcal H_{c_b}B_e)\\
\mathcal S_b\mathcal L_{c_b}B_o\\
\mathcal S_b\mathcal L_{c_b}B_e
\end{pmatrix}.
\end{equation}
Apply the complete-block update
\begin{equation}\label{eq:acceptedblock}
X\ \longleftarrow\ \sum_{j=0}^{\min(K,m_b)}
                  \binom{m_b}{j}\mathcal A_b^jX.
\end{equation}
Then $A_{\rm acc}:=\sum_iX_i(1)$ bounds the number of distinct medium
products accepted by at least one Rosser sign. The same complete-set
bound holds after arbitrary prime deletion, retaining the complete
$K,J,c_b$ and replacing only $m_b$ by the number of available primes.
\end{lemma}
\begin{proof}
At new selected length $r$, the upper sign checks odd $r$ and the lower
sign checks even $r$. If $d_{\rm old}$ is the preceding product, a checked
append of prime $p$ passes only if $d_{\rm old}p^3<D_0$.
Its preceding label $B$ has $hB\le\log d_{\rm old}$, so actual passing
implies
\[
B<\frac{\log D_0-3\log p}{h}
 \le\frac{\log D_0-3\log\min\mathcal M_b}{h},\qquad B<c_b.
\]
Consequently, a true surviving sign is never killed by the relaxed
cutoff. Unchecked appends preserve that sign. For both signs alive,
exactly one sign is checked at each append: a high-label append leaves
the other sign alone. This gives \eqref{eq:acceptedoperator}.
Each label and flag state has at most one destination; the low and high
branches are disjoint. Induction shows that the relaxed alive set
contains the true alive set, and every retained relaxed subset occurs
in exactly one of the six exclusive states. Subsets with no surviving sign may
be dropped permanently.

All primes in a block have the same label and relaxed cutoff. Choosing
$j$ distinct primes therefore applies the same operator $j$ times, with
exact multiplicity $\binom{m_b}{j}$. Their actual decreasing order does
not affect this relaxed update. Every true final accepted product is
below $H_0$, hence contains at most $K$ primes and has label below $J$.
Every preceding label is no larger, so neither the block degree cutoff
nor intermediate truncation can remove that true contribution.
Excessive total degrees and possible passes retained by rounding are
additional positive terms. The resulting sum bounds the union, without
counting a subset twice when both signs survive.

After deletion, every surviving prime still belongs to its complete
block. The same complete minimum supplies a valid necessary passing
condition; the complete $K,J$ still cover every accepted product.
The operators have nonnegative coefficients and
$\binom{m_b'}{j}\le\binom{m_b}{j}$ for $m_b'\le m_b$.
Induction over the blocks therefore gives coefficientwise domination
in every state. No subtraction of uncertain intervals is required.
\end{proof}

Replace $A_\sigma$ by $A_{\rm acc}$ in \eqref{eq:roughsupport} and set
\begin{equation}\label{eq:acceptedfraction}
\nu_{\rm acc}=\frac{2^{\omega(Q_*)}A_{\rm acc}}{Q_*H_0}.
\end{equation}
The composite coefficient support is bounded by the union of the two
medium supports and a $T$-rough large factor, so the same proof gives
$\nu_{\rm acc}(a_T+b_T/D)H$ for both sieve signs.
For the whole $A_q$ sum, absorb $q$ into the large factor as in Section
\ref{sec:roughinterval}. The medium product is still in this same union;
the unique prime at least $z$ makes $(q,d)\mapsto qd$ injective over
the entire sum. Thus the new count also bounds its aggregate image.
The integer roughness boundary, prime-modulus centering, and bilinear
charges are unchanged.

For $T=5000000$, $w_0=5$, $u=87/16$, $\kappa=11/169$ and count width
$h=1/100$, the complete count has $K=22$, $J=8115$ and 992 blocks.
It gives
\[
A_{\rm acc}=655257348020585286337311701158,
\qquad \nu_{\rm acc}<1.007455\cdot10^{-6}<1/992601.
\]
The preceding unfiltered $A_0$ is more than $8.37$ times as large.
The gap budgets still use their independently certified width $1/200$.

\paragraph{Coverage of all larger integers.}
For the lower starting point use $L_0=24274$, $L_*=24799$,
$\alpha=1057/250000$ and $\beta=93/50$; the medium family is $u=87/16$.
Equation \eqref{eq:intervalplateau}, with the new
$\nu_*:=1/992601$ and the same $a_*:=36411450583/10^{12}$,
supplies a constant leading AP expense throughout $[L_0,L_*]$.
The support condition $\alpha L-\beta\log L>\log(Q_*H_0)$,
$H\ge10^9$, $H\le\sqrt N$, $T<z$, $3<s<4$ and $k>1$ hold at
$L_0$ and persist. In particular $\alpha-\beta/L_0>0$ and
$1/2-\beta/L_0>0$. The finite charge $b_Te^{-tL_0}$ is positive
and below $10^{-5189}$; it has not been discarded.
The fixed-parameter monotonicity argument in Section \ref{sec:roughinterval}
proves the required surplus over this entire closed interval.
The preceding interval $[24799,26820]$ and the preceding half-line
$[26820,\infty)$ use the same $10^{-5}$ claim and normalization.
Together these inclusive ranges prove the first assertion of Theorem
\ref{thm:grh} for every $L\ge24274$.

For the stronger claim, use $u=11/2$, the same $T,w_0,\kappa$,
$\alpha=3737/1000000$, $\beta=3$ and $L_0=31000$.
Its gap budgets are $\eta_+=1/3231$, $\eta_-=1/3242$, with
$\epsilon=1/41219$. The accepted count gives
$\nu_{\rm acc}<8.498637\cdot10^{-7}<\nu_*:=1/1176659$, and
$a_T<a_*:=36409264976/10^{12}$.
Now $L^{3-\beta}=1$, so the leading AP expense
\[
\frac{3\nu_*}{2U_*}(a_*+b_Te^{-tL})P(L)
\]
decreases on the whole half-line. Its value at $31000$ is a valid
constant upper bound for every $L\ge31000$.
All other losses decrease and the support slack increases.
This separate half-line argument proves the stronger coefficient
$0.0181$ of Theorem \ref{thm:grh}; it does not transfer a larger
claim through an earlier certificate with a smaller surplus.
The bilinear estimate continues to use the maximum modulus $H$.

\begin{center}
\begin{tabular}{lrr}\toprule
Term & $L_0=24274$ & $L_0=31000$\\\midrule
$\alpha$ & $1057/250000$ & $3737/1000000$\\
$\beta$ & $93/50$ & $3$\\
$s$ & $3.966176$ & $3.970104$\\
$M_0$ & $>7.806728439046083$ & $>7.809030620429977$\\
$M_1$ & $<6.498811372057018$ & $<6.485781769136601$\\
$M_2$ & $<1.306485180691536$ & $<1.305068915489237$\\
$E_{\rm AP}$ & $<0.001407561467773$ & $<0.000000011604981$\\
Support slack & $>0.000548291572001$ & $>0.008509203429583$\\
Surplus & $>0.000024324829757$ & $>0.018179924199159$\\
$E_{\rm rect}$ & $<10^{-188}$ & $<10^{-250}$\\
$E_{\rm fin}$ & $<10^{-1300}$ & $<10^{-1673}$\\\bottomrule
\end{tabular}
\end{center}
The separate-formula program reproduces 12 quantities for each new
endpoint, regenerates all 348513 integer primes through $T$, and
rechecks the earlier interval and tails. Exact small-set checks cover
4608 complete subsets, 114876 integer histogram coefficients,
104976 deleted-subset patterns and 384 deleted production histograms.
The earlier 8608 aggregate prime-factor patterns remain applicable.
The 70 weak-endpoint candidates and nine stronger-count
candidates are a finite search, not a proof of global optimality.
These checks supplement the general proof above; they do not replace
the external analytic inputs or independent research review.
No distribution hypothesis beyond GRH has been introduced, and this
analytic interval join still does not reach finite verification.
% END CURRENT ACCEPTED INTERVAL
\section{A parameterized EH-style reduction}
For a quantitative comparison with standard distribution conjectures,
put
\[
\mathcal E_\theta(x)=\sum_{d\le x^\theta}\mu^2(d)
\max_{(a,d)=1}|\pi(x;d,a)-\pi(x)/\varphi(d)|.
\]
\begin{lemma}[Deleted-prime remainders]\label{lem:deleted}
If $10^9\le H\le N^\theta$, then, with the density $g(d)=1/\varphi(d)$,
\[
\sum_{\substack{d\le H\\d\mid P_N(z)}}|r_A(d)|
\le\mathcal E_\theta(N)+\om(N)\{2\sqrt N+1.1\log H\}.
\]
\end{lemma}
\begin{proof}
Let $\kappa$ be the number of deleted primes dividing $N$, and let
$\kappa_d$ count those with $d\mid N-p$. Since $(d,N)=1$,
\[
r_A(d)=\pi(N;d,N)-\frac{\pi(N)}{\varphi(d)}
-\kappa_d+\frac{\kappa}{\varphi(d)}.
\]
Each $N-p<N$ has at most $2\sqrt N$ positive divisors, and
$\kappa\le\om(N)$. Summing and using \eqref{eq:phimu} proves the claim.
\end{proof}

The usual EH conjecture provides asymptotic estimates with parameter
dependent implied constants; see \cite[Claim 2.2]{Polymath}. It does
not prescribe a numerical constant or a numerical starting point.
We retain these two quantities rather than assign them convenient values.

\begin{theorem}[Quantitative distribution inputs]\label{prop:EH}
Suppose a finite distribution input provides numbers $C>0$ and $X$
such that
\begin{equation}\label{eq:EHinput}
\mathcal E_\theta(x)\le Cx/\log^4x\quad(x\ge X).
\end{equation}
No RH or GRH is assumed. The following conclusions hold:
\[
\begin{array}{c|c|c}
\theta & \text{threshold for even }N & \pi_2(N)L^2/N\\\hline
3/4 & \max\{X,\exp(\max\{786,\sqrt{20C}\})\} & >0.01\\
9/10 & \max\{X,\exp(\max\{212,\sqrt{10C}\})\} & >0.02\\
19/20 & \max\{X,\exp(\max\{166,\sqrt{10C}\})\} & >0.02\\
99/100 & \max\{X,\exp(\max\{142,\sqrt{10C}\})\} & >0.02
\end{array}
\]
No numerical value of $C$ or $X$ is asserted by this theorem.
There are also weaker counting bounds
$\pi_2(N)>10^{-5}N/L^2$ with the same square-root constraints and
the smaller fixed starts $774$, $210$, $165$, and $141$, respectively.
\end{theorem}
\begin{proof}
Use the finite Rosser medium weights of Lemma \ref{lem:rankin},
with $D_0=T^u$, and Lemmas \ref{lem:degreewise} and
\ref{lem:prefixband} with $M=10$,
$\Sigma_E=\{0\}\cup\{a/20:2\le a\le12\}$, and
$\Gamma_E=\{0,3/25,6/25,9/25\}$.
Take the following fixed parameters:
\[
\begin{array}{c|r|r|c|c|r|r|r|c}
\theta&w_0&T&u&M&\epsilon^{-1}&\eta_+^{-1}&\eta_-^{-1}&(C_1,C_2)\\\hline
3/4&3&22000&39/8&10&1717&141&143&(113,114)\\
9/10&3&10000&35/8&10&1012&46&45&(113,114)\\
19/20&3&10000&17/4&10&1012&34&33&(113,114)\\
99/100&3&10000&17/4&10&1012&34&33&(113,114)
\end{array}
\]
For $\theta=3/4$ Lemma \ref{lem:ordered} gives
$\widehat K_T-1<0.000582284267$
and separate Rosser gaps
$\widehat R^{(\mathrm{band})}_+<0.007043548708$,
$\widehat R^{(\mathrm{band})}_-<0.006945935768$.
For $\theta=9/10$ the ordered product error is less than $0.000987573450$,
and the gaps obey
$\widehat R^{(\mathrm{band})}_+<0.021591509791$,
$\widehat R^{(\mathrm{band})}_-<0.021953570359$.
The higher two levels use the same product error and gap bounds,
less than $0.029170196992$ and $0.029947327150$, respectively.
The finite sums use all 2462 and 1227 medium primes, respectively.
They are strictly below the chosen $\epsilon,\eta_+,\eta_-$.
Take $z=N^{1/3}$, $H=N^\theta$, and
$D=H/(\kappa Q_*D_0)$, using Lemma \ref{lem:sharpsupport}.
Put $c=\log(\kappa Q_*)+u\log T$.
Then $s=3\theta-3c/L$ and $\kappa Q_*D_0D=H$.
Here $Q_*=3$ and $\kappa=5/49$ throughout; the costs are
$c<47.561$, $c<39.112$, and $c<37.961$ for the two higher levels.
At all four stated starting points, $s>2$, so $D\ge z^2$.
The refined inequalities \eqref{eq:sieveparity} apply with the
uniform product condition and Lemmas \ref{lem:parity},
\ref{lem:degreewise}, and \ref{lem:prefixband}.
Using JS Lemma 3.2 and the exact product identity gives, unconditionally,
\[
|A|V(z)\ge\frac{3U_NN}{L^2}\rho_U(L),\qquad
\rho_U(L)=1-\frac{20.508\cdot10^{-6}}{L}-2Le^{-L/3}.
\]
Indeed $z>10^{12}$ throughout, $D(z)/D_\infty\ge
\exp(-6.836\cdot10^{-6}/\log z)$, and the missing factors for
$p\mid N$ cost at most $2L/z$. The count $|A|>N/L$ is
unconditional by JS Lemma 3.1.
The normalized lower main coefficient is at least
\[
3U_*\rho_U(L)\{(1+\eta_+)(f(s)-C_2\epsilon e^{2-s})
-(\eta_++\eta_-)(2e^\gamma/s+C_1\epsilon e^{2-s})\}.
\]
For $2<s<3$, $f(s)$ increases while both subtracted corrections
decrease. Thus this coefficient increases with $L$ at every listed level.
The input costs $C/L^2$, which is at most $1/20$ or $1/10$
under the corresponding square-root constraint.
Lemma \ref{lem:deleted}, with $\om(N)\le L/\log2$, costs at most
\[
\frac{L^3}{\log2}\{2e^{-L/2}+1.1\theta Le^{-L}\},
\]
and excluding the survivor one costs $L^2e^{-L}$. Both decrease on
these ranges. The certificates give main coefficients greater than
$0.0607006692223$ at $(\theta,L)=(3/4,786)$ and
$0.1228594687806$ at $(9/10,212)$. At the higher levels they exceed
$0.1220961708146$ at $(19/20,166)$ and
$0.1354387463886$ at $(99/100,142)$.

After all charges the surpluses exceed $0.0107006692223$,
$0.0228594687806$, $0.0220961708146$, and $0.0354387463886$,
respectively.
For the weaker counting bounds keep every finite parameter and error
budget fixed, and evaluate the same expression at $L=774,210,165,141$.
The resulting full surpluses exceed $0.0003941826699$,
$0.0050969071192$, $0.0083045836547$, and $0.0166813945298$,
respectively. The same support, product, $2<s<3$, deleted-prime and
survivor-one checks hold at these starts. Their monotonicity extends
each bound to its whole half-line.
An element $N-p$ of $A$ is coprime to $N$, so it has no prime factor
omitted from $P_N(z)$. Every remaining survivor other than one therefore
has all prime factors at least $z$, and at most two prime factors,
because it is less than $z^3=N$. This proves all four conclusions.
\end{proof}

More generally, under \eqref{eq:EHinput} at the specified level,
the same argument gives a usable criterion for any fixed $2/3<\theta<1$.
Choose valid parameters of Lemma
\ref{lem:composite}, let $c=\log(Q_*H_0)$, and put
$s=3\theta-3c/L$. For $N\ge X$, $L>3\log(10^{12})$,
$T<N^{1/3}$, and $2<s<3$, define
\begin{align}\label{eq:generaldistribution}
\Xi(L)={}&3U_*\rho_U(L)
\{f(s)-C_2\epsilon e^{2-s}-\eta(2e^\gamma/s+C_1\epsilon)\}
-\frac{C}{L^2}\notag\\
&-\frac{L^3}{\log2}\{2e^{-L/2}+1.1\theta Le^{-L}\}
-L^2e^{-L}.
\end{align}
Whenever $\Xi(L)>0$, the sieve factor is positive and the preceding
proof gives $\pi_2(N)L^2/N\ge\Xi(L)>0$. All parameter and distribution
conditions must be supplied; the criterion introduces no numerical
value for $C$ or $X$. Since $s<3$, its negative corrections decrease
as $L$ increases for fixed parameters, and the positive main factor
increases. Thus a positive starting-point bound extends to the
corresponding half-line. With separate gap bounds one can replace the
sieve factor by $(1+\eta_+)l(s)-(\eta_++\eta_-)u(s)$, using
$l(s)=f(s)-C_2\epsilon e^{2-s}$ and
$u(s)=2e^\gamma/s+C_1\epsilon e^{2-s}$.
This factor is also increasing for $2<s<3$; the four rows above use
this refined version.

The existence of unspecified constants of the type in
\eqref{eq:EHinput} follows from the usual EH conjecture.
One may transfer its von Mangoldt formulation to prime counting with
a small loss in distribution level and logarithmic power, using partial
summation; ordinary EH supplies all such fixed levels below one and
all logarithmic powers. This observation does not make the constants
effective or supply $X\le e^{212}$. For instance, substituting $C=1$
and $X=e^{212}$ in the second row would be a new finite hypothesis.
We do not make that substitution or report $e^{212}$ as a numerical
threshold or record under ordinary EH alone.

\section{Signed remainders and the support actually used}
The absolute modulus sum is sufficient but is more than the lower sieve
uses. Fix one of the four parameter sets in Theorem \ref{prop:EH}.
At each even $N$ in the sieve domain use its composite lower coefficients $\lambda_N^-$,
with $z=N^{1/3}$ and $D=N^\theta/(Q_*H_0)$. Define
\[
G_N^-=\sum_d\frac{\lambda_N^-(d)}{\varphi(d)},\qquad
\mathcal R_N^-=\sum_d\lambda_N^-(d)
\left(\pi(N;d,N)-\frac{\pi(N)}{\varphi(d)}\right).
\]
These sums have precisely the support of the lower weights,
contained in $d\mid P_N(z)$, $d<N^\theta$.

\begin{lemma}[Signed lower remainder and deletions]\label{lem:signed}
For the sequence $A=\{N-p:p<N,\ p\text{ prime},\ p\nmid N\}$,
\[
S(A,z)\ge\pi(N)G_N^-+\mathcal R_N^- -\om(N).
\]
\end{lemma}
\begin{proof}
First sieve $A^*=\{N-p:p\le N,\ p\text{ prime}\}$.
Since $N$ is even and at least four, no term has $p=N$.
Summing the pointwise lower divisor inequality of Lemma
\ref{lem:composite} over $A^*$ gives
\[
S(A^*,z)\ge\sum_d\lambda_N^-(d)\pi(N;d,N)
=\pi(N)G_N^-+\mathcal R_N^-.
\]
Deleting the primes dividing $N$ removes at most $\om(N)$ survivors.
This proves the claim without taking absolute values of their divisor sums.
In particular, the $2\sqrt N$ deletion cost in Lemma \ref{lem:deleted}
is unnecessary when only this signed lower remainder is used.
\end{proof}

The following is a \emph{tailored quantitative input}, not a newly named
standard conjecture. Let $\mathfrak D_\theta$ be the even integers in
the preceding domain $L>3\log(10^{12})$, $T<z$, $2<s<3$.
For the specified family of weights, suppose
\begin{equation}\label{eq:signedinput}
\mathcal R_N^-\ge-\frac{CN}{\log^4N}
\qquad(N\ge X,\ N\in\mathfrak D_\theta).
\end{equation}
The complete input \eqref{eq:EHinput} implies \eqref{eq:signedinput},
with the same $C,X$, because $|\lambda_N^-(d)|\le1$.
Only one side of one signed sum is required here. Neither GRH nor a
published weighted asymptotic estimate supplies its numerical $C,X$.
The integrality lower bound for the complete maximum-error sum does not,
by itself, give a lower bound for this signed sum.

Put $c=\log(Q_*H_0)$, $s=3\theta-3c/L$, and use the separate gaps of
Theorem \ref{prop:EH}. With its functions $l(s),u(s)$, write
\[
M_\theta(L)=3U_*\rho_U(L)
\{(1+\eta_+)l(s)-(\eta_++\eta_-)u(s)\}.
\]
Lemma \ref{lem:signed}, the same product and prime-count inputs, and
$\om(N)\le L/\log2$ give the sharper criterion
\begin{equation}\label{eq:signedcriterion}
\pi_2(N)\frac{L^2}{N}\ge
\Xi_{\rm w}(L):=M_\theta(L)-\frac{C}{L^2}
-\left(\frac{L^3}{\log2}+L^2\right)e^{-L},
\end{equation}
when the main factor is positive and the preceding sieve domains hold.
The last $L^2e^{-L}$ charge removes the possible survivor one.
For a requested surplus $a>0$, the admissible budget at $L$ is therefore
\[
C<L^2\left\{M_\theta(L)
-\left(\frac{L^3}{\log2}+L^2\right)e^{-L}-a\right\}.
\]
This spends the actual positive main mass, instead of imposing the
fixed $1/20$ or $1/10$ budgets of Theorem \ref{prop:EH}.

\begin{proposition}[Finite-constant budget conversion]\label{prop:weighted}
Suppose \eqref{eq:signedinput} holds for one of the four specified weight
families. Each entry $A$ below gives
$\pi_2(N)>10^{-5}N/\log^2N$ for every even $N\ge\max\{X,e^A\}$,
provided $C$ does not exceed the indicated column label:
\[
\begin{array}{c|rrrr}
\theta&C\le1&C\le10^3&C\le10^4&C\le10^6\\\hline
3/4&722&723&740&1632\\
9/10&200&202&222&936\\
19/20&158&161&182&872\\
99/100&136&138&159&829
\end{array}
\]
The same conclusions hold under the complete input \eqref{eq:EHinput}.
The column labels are allowed budgets, not distribution constants proved
for the primes. The onset $X$ is retained in every conclusion.
\end{proposition}
\begin{proof}
The 192-bit certificate evaluates \eqref{eq:signedcriterion} strictly
above $10^{-5}$ at each entry, checks all sieve domains, and reproduces
the saved main masses of Theorem \ref{prop:EH} within $10^{-45}$.
It also evaluates the main mass by an independent algebraic expression
in which the $e^\gamma$ factors are cancelled explicitly.
On $2<s<3$, $l(s)$ increases and $u(s)$ decreases, so the composite
factor increases. The positive $\rho_U(L)$ increases on these ranges.
Both exponential charges and $C/L^2$ decrease. Hence the inequality
extends to the whole half-line. For smaller $C$ the bound only improves.
The $P_2$ transfer is the same as in Theorem \ref{prop:EH}.
\end{proof}
For example, if a quantitative input is supplied with upper constant
$C=10^6$ at level $99/100$, the earlier square-root allocation requires
$L\ge\sqrt{10^7}=3162.277\ldots$, whereas Proposition
\ref{prop:weighted} requires $L\ge829$. Both retain the same unknown
onset $X$. This is an improvement of the budget conversion under the
same input; it is not a numerical record under ordinary EH.
Even with zero charged remainder the current certified expression has
negative main factor at $L=721,199,157,135$ for the four rows.
Its first positive integer starts are $722,200,158,136$.
This limitation concerns these fixed weights and main-mass bounds.

% BEGIN CURRENT PREFIX DISTRIBUTION
\subsection{Prefix-block finite-constant conversion}\label{sec:prefixdistribution}
Recompute the finite medium masses using Lemma \ref{lem:prefixblocks}.
The following five families use $w_0=3$, $Q_*=3$, $\kappa=5/49$,
$D_0=T^u$, and the indicated gap and product budgets:
\begin{center}
\small
\begin{tabular}{rrrrrrrr}\toprule
Family & $T$ & $u$ & $h$ & $\epsilon$ & $\eta_+$ & $\eta_-$ & $(C_1,C_2)$\\\midrule
$1$ & $22000$ & $33/8$ & $1/200$ & $1/1717$ & $1/100$ & $1/103$ & $(113,114)$\\
$2$ & $100000$ & $9/2$ & $1/100$ & $1/4729$ & $1/233$ & $1/235$ & $(109,110)$\\
$3$ & $10000$ & $29/8$ & $1/200$ & $1/1012$ & $1/32$ & $1/31$ & $(113,114)$\\
$4$ & $10000$ & $15/4$ & $1/200$ & $1/1012$ & $1/42$ & $1/42$ & $(113,114)$\\
$5$ & $10000$ & $7/2$ & $1/200$ & $1/1012$ & $1/24$ & $1/23$ & $(113,114)$\\
\bottomrule\end{tabular}
\end{center}
For a chosen cell, use that family's composite lower weights at
$z=N^{1/3}$ and $D=N^\theta/(\kappa Q_*D_0)$.
If only the signed input \eqref{eq:signedinput} is assumed, it must hold
for this specified family on its sieve domain. The full absolute input
\eqref{eq:EHinput} implies it for every family with the same $C,X$.
No value of $C$ or $X$ is supplied by this construction.

\begin{proposition}[Prefix-block budget conversion]\label{prop:prefixdistribution}
Under that cell's signed input, or the full absolute input, every even
$N\ge\max\{X,e^A\}$ satisfies $\pi_2(N)>10^{-5}N/\log^2N$.
Here each entry is $A$ followed by its family number, and $C$ is bounded
by the column label:
\[
\begin{array}{c|rrrr}
\theta&C\le1&C\le10^3&C\le10^4&C\le10^6\\\hline
$3/4$ & $627\ (1)$ & $629\ (1)$ & $648\ (1)$ & $1510\ (2)$\\
$9/10$ & $171\ (3)$ & $174\ (3)$ & $198\ (4)$ & $885\ (2)$\\
$19/20$ & $135\ (5)$ & $138\ (5)$ & $162\ (3)$ & $825\ (2)$\\
$99/100$ & $115\ (5)$ & $118\ (5)$ & $143\ (3)$ & $791\ (2)$\\
\end{array}
\]
\end{proposition}
\begin{proof}
The complete-set bounds of Lemma \ref{lem:prefixblocks} supply the two
gaps in \eqref{eq:sieveparity}; all coefficient, support and signed
remainder identities continue to apply. Use exactly
\eqref{eq:signedcriterion} with each family's $c=\log(\kappa Q_*)+u\log T$.
The 192-bit certificates check every cell's strict surplus, $T<z$,
$L>3\log(10^{12})$, and $2<s=3\theta-3c/L<3$.
For fixed parameters, $l(s)$ increases and $u(s)$ decreases on $(2,3)$,
$s$ and $\rho_U(L)$ increase, and every deducted expense decreases.
Thus every cell extends to the claimed half-line.
The exact and separate-formula checks test finite arithmetic, not a
distribution upper bound or independent mathematical review.
\end{proof}
Ordinary EH still supplies only unspecified constants and starting points.
These entries improve the conversion of a supplied quantitative input,
and are not numeric EH-only records. In particular $X$ cannot be discarded.
The smallest support cost among these five families is 31.052421204;
even $\theta<1$ allows only $L_B/3<14.278$ at $B$.
Thus these families still cannot reach the verified endpoint, even with
zero signed error. The complete absolute input also retains the integrality
constraints of Proposition \ref{prop:tradeoff}.
% END CURRENT PREFIX DISTRIBUTION
\begin{proposition}[A support obstruction to triple factorability]
\label{prop:factorability}
The present four lower-weight families are not triply well-factorable
of level $N^\theta$ on a whole half-line.
\end{proposition}
\begin{proof}
For a prime $T<p<z$ with $p\nmid N$, the lower Rosser weight has no
checked even prefix, and the empty medium gap is zero. Consequently
$\lambda_N^-(p)=-1$. A triply well-factorable sequence of level $H$
must have coefficient zero at every prime $p>H^{1/3}$: apply the
balanced split $H_1=H_2=H_3=H^{1/3}$ in the definition
\cite[Definition 1.1]{Pascadi}. Every factorization of a prime has a
factor equal to $p$, outside its allowed support.

For the four rational levels, take $N=2^{3000k}$, $k\ge1$.
Then $H^{1/3}=2^{1000\theta k}$ is an integer. Bertrand's theorem
provides a prime strictly between $H^{1/3}$ and $2H^{1/3}<z$.
It exceeds $T$ and cannot divide $N$. The parameter certificate
checks $D\ge z^2$ already at $k=1$, and that support slack increases
with $k$. This prime has coefficient $-1$, a contradiction.
The certificate also supplies the exact small witness
$N=10^{18}$, $z=10^6$, $T=5$, $D_0=625$, $D=10^{14}$:
$p=999983$ is prime, its coefficient is $-1$, and
$p^3>H=3D_0D=187500000000000000$.
\end{proof}
This does not decide two-factor well-factorability or the applicability
of a separately constructed well-factorable linear-sieve variant.
Changing the coefficient family requires a new pointwise and main-mass
proof; support size alone is not enough.

Pascadi's fixed-residue theorem gives $5/8-\varepsilon$ for triply
well-factorable weights \emph{unconditionally}, and $3/5-\varepsilon$
for its upper linear-sieve variant; Theorem 5.4 treats specified lower
blocks with factor-dependent levels \cite{Pascadi}. Its constants
depend on the residue $a$, so setting $a=N$ does not give the required
uniform input. Lichtman's large-residue formula is
$(5-\vartheta_{\rm RP})/8-\varepsilon$, giving $153/256-\varepsilon$
at $\vartheta_{\rm RP}=7/32$ \cite[Theorem 1.7]{Lichtman}.
Setting $\vartheta_{\rm RP}=0$ requires the Ramanujan--Petersson
input identified there, rather than Dirichlet GRH alone.
Even the resulting $5/8$ is below the $2/3$ needed for this direct
$z=N^{1/3}$ linear lower sieve; Chen switching remains necessary
for such lower distribution levels.

The MPZ smooth-modulus input allows a primitive residue depending
on $N$, but requires prime factors of its moduli at most $N^\delta$.
The proved region in \cite[Claim 2.4, Theorem 2.5]{Polymath} is
$600\varpi+180\delta<7$ at level $1/2+2\varpi$.
The current nonzero prime coefficients extend to $N^{1/3}$ on a
half-line, so their uniform coverage needs $\delta\ge1/3$:
apply Bertrand near $z$ along $N=2^{3000k}$.
This is outside that region, whose level supremum is only
$157/300<2/3$. A restriction to smooth moduli would leave
uncontrolled coefficients and needs a different decomposition.

Li's 2026 theorem states the asymptotic lower count
$1.9728C_NN/\log^2N$, combining weighted sieves, switching, the
double sieve and Harman's sieve \cite{Li2026}.
Its statement does not supply a numerical starting point.
Explicit uniform estimates for its switched blocks are a route
to investigate, not an explicit threshold certified here.

\section{Cost calculation for the rejected bridge candidate}
\textbf{Withdrawal notice.} Proposition \ref{prop:integrality} refutes
the distribution premise of this section. The following algebra is
retained to identify the original error: it checks that a proposed
budget would suffice, without checking whether that budget is possible.
It does not prove a nonvacuous full-range result.
Here we do not use the GRH rectangle machinery. Put $z=N^{1/3}$,
$H=N^{9/10}$, $D=H/e^8$, $N\ge B$. Lemma \ref{lem:product}
permits $\epsilon=1/700$, $C_2=117$, and
\[
s(L)=2.7-24/L\ge s(L_B)=2.139682079805\ldots>2.
\]
The former calculation had $D\ge z^2$ and $Q_ND\le H$; its
fixed-endpoint-only sieve application was not independently justified.
For $2\le s\le2.7$, $f(s)$ increases while $e^2h(s)=e^{2-s}$
decreases. Hence the positive sieve factor is at least
\[
f_B=f(s(L_B))-(117/700)e^{2-s(L_B)}
>0.0723184176500.
\]
Also \eqref{eq:Vbounds} and \eqref{eq:A} give
\[
|A|V(z)\ge\frac{3U_NN}{L^2}\rho(L),\qquad
\rho(L)=1-\frac{L}{8\pi}e^{-L/6}-2Le^{-L/3}.
\]
The function $\rho$ increases on this range, and
$\rho(L_B)>0.9985933587896$.

QEH$_B$ concerns all primes, while $A$ excludes $p\mid N$.
We explicitly account for the difference. Write
$\kappa=\#\{p<N:p\mid N\}$ and
$\kappa_d=\#\{p<N:p\mid N,\ d\mid N-p\}$.
For $(d,N)=1$,
\[
r_A(d)=\pi(N;d,N)-\frac{\pi(N)}{\varphi(d)}
-\kappa_d+\frac{\kappa}{\varphi(d)}.
\]
For each deleted prime $p$, the integer $N-p$ has at most
$2\sqrt N$ positive divisors. Therefore, irrespective of $H$,
\begin{equation}\label{eq:del}
\sum_{\substack{d\le H\\d\mid P(z)}}|r_A(d)|
\le\frac{N}{10L^2}
+\om(N)\{2\sqrt N+1.1\log H\}.
\end{equation}
This divisor count avoids the wasteful $\om(N)H$ estimate.
Use Robin's bound $\om(N)<1.3841L/\log L$, recorded in \cite{BS}.
After normalizing by $N/L^2$ and excluding the possible survivor one,
\begin{align}\label{eq:bridgemargin}
\pi_2(N)\frac{L^2}{N}\ge{}&3e^{-\gamma}\rho(L_B)f_B-\frac1{10}
\\&-\frac{1.3841L^3}{\log L}
\left(2e^{-L/2}+0.99Le^{-L}\right)-L^2e^{-L}.\notag
\end{align}
An element of $A$ surviving $P_N(z)$ is coprime to $N$, hence has
no excluded prime factor. Its prime factors are all at least $z$.
Three factors would give an integer at least $z^3=N$, whereas every
element of $A$ is less than $N$. So every surviving element other
than one is a legitimate prime or semiprime.
The error terms on the second line of \eqref{eq:bridgemargin} decrease
for $L\ge L_B$. At $B$ the main coefficient exceeds
$0.1216402390842$; the deleted-prime budget is less than
$0.0000289481782$ and $L_B^2/B<4.59\cdot10^{-16}$.
The final surplus exceeds $0.0216112909060>0.021$.
The computation of \cite{OHP} covers every even $4\le N\le B$
by two primes. Nevertheless the proposed analytic premise fails
at $B$, so this calculation does not establish a bridge.

\section{Interpretation and remaining research questions}
The GRH theorem changes finite constants and presieving cost,
while retaining the asymptotic square-root distribution level.
The composite weights resolve the endpoint problem for this theorem,
but its threshold still does not bridge the verified range.
Under the usual GRH absolute-error summation,
direct sieving at $z=N^{1/3}$ has asymptotic parameter $s\le3/2$;
the lower linear sieve has no positive main term there.
Chen's switching avoids that obstruction at large $N$, with a
finite support and error cost that remains large near $B$.
For the present direct distribution route, $D=N^\theta/e^c$ and
$z=N^{1/3}$ require $c\le(\theta-2/3)L$ just to have $D\ge z^2$.
At $B=4\cdot10^{18}$, $L_B=42.832826035\ldots$, the allowances are
only $3.569402170$ and $9.994326075$ at the two chosen levels, whereas
the earlier fixed-family costs are $47.560368848\ldots$ and $39.111469030\ldots$.
Thus even zero AP error would not make these fixed weights reach $B$.
A substantially smaller support cost or a different sieve construction
is needed for this route. This is an obstruction to the present
construction, not an impossibility theorem for conditional Chen bounds.

The rejected QEH$_B$ candidate changes the level and supplies an
explicit constant and starting point. Its sieve budget would have
been sufficient, but its distribution budget is impossible at $B$.
Any replacement must first pass the integrality obstruction and then
be related to a recognized conjecture or independently supported finite
distribution input. Whole-range stitching remains open in this project.
RH alone controls complete prime products and global prime counts,
and has not supplied the AP input required by Theorem \ref{thm:grh}.
Finite-height GRH verification cannot be substituted for full GRH
without a certified explicit-formula tail.

For whole-range comparisons the 2026 thesis summary of Johnston
\cite{Jthesis} states GRH $p+P_{16}$, improving the older $p+P_{31}$.
It must not be confused with a whole-range $p+P_2$ theorem.
The withdrawn candidate supplies no improvement on this whole-range result.
EH and GEH are recognized conjectures used for conditional prime-gap
results \cite{Polymath}. GRH is not known to imply either. Their usual
forms do not provide the concrete constants and starting points needed
to connect a finite verification with an explicit analytic threshold.
Theorem \ref{prop:EH} records this dependence instead of hiding it.

Proposition \ref{prop:tradeoff} also constrains the optimization itself:
the level $99/100$ route of Theorem \ref{prop:EH} has an actual threshold
greater than $e^{834}$ under that theorem's square-root budget.
Merely relaxing that budget cannot recover
an endpoint near $e^{141}$ within the present direct-sieve expression.
Indeed, for \eqref{eq:generaldistribution} and its separate-gap
version, assume the proved conditions $0<\rho_U(L)\le1$,
$2<s\le3\theta<3$, nonnegative support cost, and nonnegative error
parameters. Since $f(s)<F(s)$ and $f$ increases on this interval,
the main coefficient is at most
\[
3U_*f(3\theta)
=\frac{7.92\log(3\theta-1)}{3\theta}
=:K_\theta.
\]
All other charges are nonnegative, so $\Xi(R)>0$ requires
$C<K_\theta R^2$. At $\theta=99/100$, the strict certificate gives
$\mathfrak I_\theta(492)>K_\theta\,492^2$.
If such a positive endpoint satisfied $R\le492$ and $e^R\ge X$,
the input at $x=e^{492}$ would imply
$C>\mathfrak I_\theta(492)>K_\theta\,492^2\ge K_\theta R^2$,
a contradiction. Thus even after reducing support cost and choosing
a larger error budget, this direct-sieve expression needs $R>492$.
This conclusion uses the stated normalization and main-term model;
it is not a bound for all conditional Chen arguments.

\section*{Reproducibility and dependency ledger}
% BEGIN ACCEPTED INTERVAL LEDGER
For the accepted-support union and current endpoints, run:
\begin{quote}\small
\path{conditional_chen_accepted_support_checks.py}\\
\path{conditional_chen_accepted_interval_certificate.py}\\
\path{conditional_chen_accepted_interval_checks.py}
\end{quote}
The counter uses six positive exclusive flag/parity states and exact
binomial block multiplicities. Its complete cutoffs remain fixed after
deletion. Counts are generated completely or reused only after generator
hash checks. The scalar checker uses its own endpoint formulas and
integer prime sieve, checks exact sums of the six state counts, and
rechecks the preceding interval and tails. The weak interval covers
$[24274,24799]$ and joins the earlier $[24799,26820]$ and GRH tail;
the stronger coefficient uses its own fixed-$\beta=3$ half-line.
Continuous coverage follows from Section \ref{sec:acceptedinterval},
not from finitely many sample points. Neither a finite-verification
splice nor independent research review is supplied by these programs.
% END ACCEPTED INTERVAL LEDGER
% BEGIN ROUGH INTERVAL LEDGER
For the logarithmic count and rough-integer interval, run:
\begin{quote}\small
\path{conditional_chen_binned_support_checks.py}\\
\path{conditional_chen_rough_interval_certificate.py}\\
\path{conditional_chen_rough_interval_checks.py}
\end{quote}
The support histograms are generated completely or reused from per-case
caches after checking generator hashes. The interval program computes the
complete integer Euler product freshly; the separate checker regenerates
its prime set independently. It rechecks the old tail, verifies the common
claim and inclusive join, and uses its own formulas for 12 new endpoint
quantities. Continuous coverage follows from the fixed-parameter argument
in Section \ref{sec:roughinterval}, not from the five sample points.
This work does not supply independent research review or a splice to the
finite verification range.
% END ROUGH INTERVAL LEDGER
For the current prefix-block results, run the following files in order:
\begin{quote}\small
\path{conditional_chen_prefix_block_checks.py}\\
\path{conditional_chen_prefix_block_endpoint_certificate.py}\\
\path{conditional_chen_prefix_block_refinement.py}
\end{quote}
The refinement records
per-case caches tied to generator hashes. The initial grid includes a
previously generated $u=11/2$ finite sum, explicitly reused by hash.
The separate support-count program
\path{conditional_chen_sparse_modulus_certificate.py} freshly evaluates
the complete Euler products at four rational exponents for each of three
cutoffs; it also checks 18 exact small Rankin counts.
Then run \path{conditional_chen_prefix_block_endpoint_checks.py}.
It independently reproduces eight GRH quantities at each asserted endpoint
and the main mass and surplus in all 16 distribution cells, without calling
the production scalar evaluators. The exact prefix-block checks cover
4608 deleted prime subsets, 9216 first-failure identities, 104976 divisor
patterns, and 132 block-coefficient groups. These tests do not replace
Lemmas \ref{lem:prefixblocks} and \ref{lem:sparsecount}, the external inputs,
independent research review, or a verification-range splice.

First run \texttt{python}\ \path{conditional_chen_composite_certificate.py}
to regenerate the finite inputs and the conservative-support baseline.
Then run \texttt{python}\ \path{conditional_chen_sharp_endpoint_certificate.py}
for the preceding GRH endpoints and the distribution theorem.
For the preceding dense-grid GRH endpoints, then run
\texttt{python}\ \path{conditional_chen_dense_band_grh_certificate.py}.
It freshly generates the enlarged-grid finite sums, saves
\path{conditional_chen_dense_band_grh_finite.json}, and checks the
new AP budget and endpoints. The optional argument
\texttt{--reuse-finite} reuses that file after checking its generator
hashes. Five quantities at the earlier $L=30824$ are reproduced
within $10^{-45}$; direct AP, surplus and support formulas are also
checked within $10^{-45}$. The ordered product input is reused with
its hash, not regenerated in this additional run.
The sharp endpoint program records the SHA-256 hashes of
the saved finite inputs, scans the complete medium-prime sets again for
the support maxima, and checks all new support, parameter and surplus
conditions. Its independent scalar implementation also reproduces eight
baseline GRH quantities within $10^{-45}$ before recharging the support.
The first program regenerates the
complete ordered-product certificate through $10^8$ by calling
\path{conditional_chen_uniform_product_certificate.py}; saved
JSON values are not used as proof inputs in that first program.
The independent finite
\path{conditional_chen_uniform_product_checks.py} checks 4096
integer sequences against exhaustive ordered pairs, 1954 exact prime
product endpoint pairs, and nine integer/prime tail comparisons.
Those checks exercise the merging and endpoint conventions, and do
not replace Lemma \ref{lem:ordered} or its published tail inputs.
The module
\path{conditional_chen_rankin_certificate.py} generates all
medium primes by a complete integer sieve and sums the positive Rankin
bound in 192-bit Arb. The earlier
\path{conditional_chen_degreewise_rankin_certificate.py}
uses its integer sieve and degree thresholds, and maintains the positive
recurrence separately for each exponent. The earlier
\path{conditional_chen_quadratic_rankin_certificate.py}
adds the moments at $0,\sigma,2\sigma$, checks the optimizer signs
strictly, and encloses the nonnegative quadratic expression.
It is retained as a baseline and can also be run separately.
The separate run writes a \texttt{\_standalone.json} file, preserving
the production JSON with its freshly computed uniform-product metadata.
The separate \path{conditional_chen_degreewise_rankin_checks.py}
compares nine finite coefficient configurations with exact rational
subset sums and verifies 2304 first-failure and prime-deletion cases.
The independent \path{conditional_chen_quadratic_rankin_checks.py}
also checks nine configurations and 2304 cases, evaluating the optimized
polynomial by exact rational subset sums and comparing actual first
failures for all deleted prime subsets.
The current \path{conditional_chen_prefix_band_certificate.py}
maintains positive raw exponential-generating moments through degree four,
encloses the signed shift and Gram optimization by Arb, and skips any
optimizer whose signs cannot be certified. Native truncated polynomial
multiplication computes the required coefficients exactly, without omitting
any subset mass. Its separate run preserves the production JSON.
The independent \path{conditional_chen_prefix_band_checks.py}
uses formal power-of-two atoms, not primes, so that base-two logarithms
and polynomial expectations are exact rational numbers. It checks twelve
configurations, 2688 first-failure/deletion cases, and 25 actual product-band
instances, including one two-atom first failure. These checks validate the
finite algebra; they do not replace the general prime-set proof.
The \path{conditional_chen_prefix_band_eh_search.py} explores 96
strict scalar candidates. Its four selected finite inputs are regenerated
by the composite certificate; the sharp endpoint program then selects the
new starts with the smaller proved support costs. The earlier quadratic
search of 80 and degreewise search of 144 candidates remain baselines.
The \path{conditional_chen_sharp_support_certificate.py}
computes the exact support maxima by complete adjacent-prime scans and
checks 52488 exact coefficient/deletion cases.
The older
\path{conditional_chen_product_certificate.py} and
\path{conditional_chen_interval_certificate.py} are retained
for the unused fixed-endpoint candidates and withdrawn bridge costs.
The separate
\path{conditional_chen_integrality_certificate.py} refutes both
finite EH-style candidates. Only strict Arb comparisons accept a claim.
The \path{conditional_chen_distribution_tradeoff_certificate.py}
certifies the necessary constant/onset witnesses in Proposition
\ref{prop:tradeoff} and the direct-sieve positivity obstruction above.
It also computes actual reduced-class prime histograms for 19701
integer pairs $3\le x\le200$, $2\le H<x$, checking the prime-modulus
subtotal against the integrality lower bound with exact rational
counts. These finite checks do not replace the general proof or
provide any distribution upper bound.
The \path{conditional_chen_weighted_distribution_certificate.py}
reproduces the saved distribution main masses and certifies Proposition
\ref{prop:weighted}. It uses exact rational arithmetic in 10748 composite
cases over 1156 selected even integers, checks 99158 divisor patterns and
21496 signed identities, including deleted primes and the $P_2$ transfer.
It verifies the finite prime witness in Proposition \ref{prop:factorability}
and the support inequalities for its infinite families.
It does not prove a signed distribution upper bound, regenerate the
large finite prime products, or certify global optimality.
The exact-integer script \path{conditional_chen_weight_checks.py}
also checks 2347008 pointwise divisor patterns across six finite
exact/Brun/Rosser configurations, including the integer levels on either side
of every weight transition. These finite checks exercise the coefficient,
support and pointwise inequalities; they do not replace the general
proof of Lemma \ref{lem:composite} or the external main-mass estimates.
The \path{conditional_chen_rosser_checks.py} checks 602368
pointwise patterns across three finite exact/Rosser/Rosser configurations,
eight additional rational-power configurations, 16 normalized coefficient
cases, and 3088 exact first-failure and prime-deletion cases. The integer
configurations include both sides of every weight transition; the rational-power
ones compare levels by exact integer powers, without numerical roots.
They also check the separate parity bounds by explicit subset sums and
the minimum-prime degree restrictions by integer powers, including the
eight rational powers $5$, $9/2$, $35/8$, $25/4$, $17/4$,
$49/8$, $39/8$, and $33/8$.
These checks do not replace Lemmas \ref{lem:rankin}, \ref{lem:parity},
\ref{lem:degreewise}, \ref{lem:quadratic}, and \ref{lem:prefixband}.
The script \path{conditional_chen_source_audit.py} checks source
structure, references, and agreement of the current manuscript and notes
with the saved certificates. It records input hashes for reproducibility;
it does not establish the analytic dependencies or compile the manuscript.
Neither the scripts nor successful compilation mechanically verify the
analytic lemmas. The dependencies to check in a research review are:
\begin{enumerate}
\item The Rosser construction and main-mass estimates in BJS
Proposition 13, Table 1, and Nathanson Theorem 9.3. The uniform product
condition in Lemma \ref{lem:uniform} is verified at every endpoint.
Lemma \ref{lem:ordered} tightens that condition by a complete finite
computation and an unconditional analytic tail.
Lemma \ref{lem:rankin} derives the finite medium-prime mass identity
directly by first failure, and Lemma \ref{lem:composite} proves the
presieving step and its coefficient and support bounds.
Lemma \ref{lem:degreewise} permits independent exponents for disjoint
degree groups. Lemma \ref{lem:quadratic} uses a nonnegative polynomial
majorant and finite weighted Cauchy--Schwarz; optimization follows
complete-prime domination. Lemma \ref{lem:prefixband} retains the
preceding-pass constraint and uses a nonnegative tilted square with a
finite weighted Gram matrix. None changes the weights.
Lemma \ref{lem:sharpsupport} sharpens their support bound without changing
the coefficients, their defining level, or their finite gap estimates.
\item The AP and bilinear estimates in the accessible BS preprint.
Their constants are used directly; its specialized three-term
threshold formulas and normalization are not used.
\item BJS Lemma 16, JS Lemma 3.2, Lee--Nosal Corollary 1.3, the square-free reciprocal
totient bound, Robin's $\om$ bound, and Nathanson's combinatorial inequality.
\item Theorem \ref{prop:EH} retains the unspecified $C,X$ in
\eqref{eq:EHinput}; no effective EH constants have been supplied.
\item Proposition \ref{prop:weighted} retains the same missing numerical
data, with a weaker signed input proved sufficient in Lemma \ref{lem:signed}.
Its smaller budgets do not supply that input or bypass the product and
support conditions. The cited weighted distribution results require
their stated coefficient classes and residue uniformity.
\item The finite distribution hypothesis in \eqref{eq:QEH} is refuted
by Proposition \ref{prop:integrality}; there is no full-bridge theorem.
\end{enumerate}
The old floating-point probe is exploratory and is superseded by this
ledger. No new unproved minorant module is
hidden inside the retained cost calculation, but its distribution
premise is impossible. This distinguishes positive cost arithmetic
from a viable conditional research result.

\begin{thebibliography}{99}
\bibitem{Pascadi} A. Pascadi,
\emph{On the exponents of distribution of primes and smooth numbers},
2025, version 2. \url{https://arxiv.org/abs/2505.00653v2}.
\bibitem{Lichtman} J. D. Lichtman,
\emph{Primes in arithmetic progressions to large moduli, and Goldbach
beyond the square-root barrier}.
\url{https://arxiv.org/abs/2309.08522}.
\bibitem{Li2026} R. Li,
\emph{On Chen's theorem, Goldbach's conjecture and almost prime twins II},
Math. Reports 28(78) (2026), 39--61.
\url{https://imar.ro/journals/Mathematical_Reports/Pdfs/2026/1_2/3.pdf}.
\bibitem{BJS} M. Bordignon, D. R. Johnston, V. V. Starichkova,
\emph{An explicit version of Chen's theorem and the linear sieve},
Int. J. Number Theory 21 (2025), 2497--2572.
\url{https://arxiv.org/abs/2207.09452v6}.
\bibitem{BS} M. Bordignon, V. Starichkova,
\emph{An explicit version of Chen's theorem assuming the Generalized
Riemann Hypothesis}, accessible preprint, 2022.
\url{https://arxiv.org/abs/2211.08844}.
\bibitem{BSpub} Same authors, published version, Ramanujan J. 64
(2024), 1213--1242. The publisher's abstract states $e^{e^{14}}$.
\url{https://doi.org/10.1007/s11139-024-00866-x}.
\bibitem{JS} D. R. Johnston, V. V. Starichkova,
\emph{Some explicit results on the sum of a prime and an almost prime}.
\url{https://arxiv.org/abs/2208.01229}.
\bibitem{LN} E. S. Lee, P. Nosal,
\emph{Sharper bounds for the error in the prime number theorem assuming
the Riemann Hypothesis}, version 4 (2025).
\url{https://arxiv.org/abs/2312.05628v4}.
\bibitem{DJ} A. W. Dudek, D. R. Johnston,
\emph{Almost primes between all squares}.
\url{https://arxiv.org/abs/2501.18048}.
\bibitem{Jthesis} D. R. Johnston,
\emph{New explicit and asymptotic results in sieve theory}, thesis
summary, Bull. Aust. Math. Soc. 113 (2026), 372--374.
\url{https://doi.org/10.1017/S0004972725100828}.
\bibitem{OHP} T. Oliveira e Silva, S. Herzog, S. Pardi,
\emph{Empirical verification of the even Goldbach conjecture and
computation of prime gaps up to $4\cdot10^{18}$}, Math. Comp. 83
(2014), 2033--2060.
\url{https://doi.org/10.1090/S0025-5718-2013-02787-1}.
\bibitem{Polymath} D. H. J. Polymath,
\emph{Variants of the Selberg sieve, and bounded intervals containing
many primes}, Research in the Mathematical Sciences 1 (2014), 12.
\url{https://arxiv.org/abs/1407.4897}.
\bibitem{Nathanson} M. B. Nathanson,
\emph{Additive Number Theory: The Classical Bases}, Graduate Texts
in Mathematics 164, Springer, 1996.
\url{https://doi.org/10.1007/978-1-4757-3845-2}.
\end{thebibliography}
\end{document}
